%\VignetteEncoding{UTF-8}
%\SweaveOpts{concordance=TRUE}
\makeatletter
\def\input@path{{./}{vignettes/manuscript/}}
\makeatother
\documentclass[article]{jss}\usepackage[]{graphicx}\usepackage[]{xcolor}
% maxwidth is the original width if it is less than linewidth
% otherwise use linewidth (to make sure the graphics do not exceed the margin)
\makeatletter
\def\maxwidth{ %
  \ifdim\Gin@nat@width>\linewidth
    \linewidth
  \else
    \Gin@nat@width
  \fi
}
\makeatother

\definecolor{fgcolor}{rgb}{0.345, 0.345, 0.345}
\newcommand{\hlnum}[1]{\textcolor[rgb]{0.686,0.059,0.569}{#1}}%
\newcommand{\hlsng}[1]{\textcolor[rgb]{0.192,0.494,0.8}{#1}}%
\newcommand{\hlcom}[1]{\textcolor[rgb]{0.678,0.584,0.686}{\textit{#1}}}%
\newcommand{\hlopt}[1]{\textcolor[rgb]{0,0,0}{#1}}%
\newcommand{\hldef}[1]{\textcolor[rgb]{0.345,0.345,0.345}{#1}}%
\newcommand{\hlkwa}[1]{\textcolor[rgb]{0.161,0.373,0.58}{\textbf{#1}}}%
\newcommand{\hlkwb}[1]{\textcolor[rgb]{0.69,0.353,0.396}{#1}}%
\newcommand{\hlkwc}[1]{\textcolor[rgb]{0.333,0.667,0.333}{#1}}%
\newcommand{\hlkwd}[1]{\textcolor[rgb]{0.737,0.353,0.396}{\textbf{#1}}}%
\let\hlipl\hlkwb

\usepackage{framed}
\makeatletter
\newenvironment{kframe}{%
 \def\at@end@of@kframe{}%
 \ifinner\ifhmode%
  \def\at@end@of@kframe{\end{minipage}}%
  \begin{minipage}{\columnwidth}%
 \fi\fi%
 \def\FrameCommand##1{\hskip\@totalleftmargin \hskip-\fboxsep
 \colorbox{shadecolor}{##1}\hskip-\fboxsep
     % There is no \\@totalrightmargin, so:
     \hskip-\linewidth \hskip-\@totalleftmargin \hskip\columnwidth}%
 \MakeFramed {\advance\hsize-\width
   \@totalleftmargin\z@ \linewidth\hsize
   \@setminipage}}%
 {\par\unskip\endMakeFramed%
 \at@end@of@kframe}
\makeatother

\definecolor{shadecolor}{rgb}{.97, .97, .97}
\definecolor{messagecolor}{rgb}{0, 0, 0}
\definecolor{warningcolor}{rgb}{1, 0, 1}
\definecolor{errorcolor}{rgb}{1, 0, 0}
\newenvironment{knitrout}{}{} % an empty environment to be redefined in TeX

\usepackage{alltt}

%% ---- Packages ----
\usepackage[utf8]{inputenc}
\usepackage[T1]{fontenc}
\usepackage{amsmath,amssymb}
\usepackage{booktabs}
\usepackage{graphicx}
\usepackage{natbib}
\usepackage{xurl}
\makeatletter
\@ifundefined{knitrout}{\newenvironment{knitrout}{}{}}{}
\makeatother
\providecolor{fgcolor}{rgb}{0,0,0}
\providecolor{bgcolor}{rgb}{1,1,1}
\providecolor{shadecolor}{rgb}{0.97,0.97,0.97}
\providecolor{messagecolor}{rgb}{0,0,0}
\providecolor{warningcolor}{rgb}{1,0,1}
\providecolor{errorcolor}{rgb}{1,0,0}
\makeatletter
\def\maxwidth{\ifdim\Gin@nat@width>\linewidth\linewidth\else\Gin@nat@width\fi}
\def\maxheight{\ifdim\Gin@nat@height>\textheight\textheight\else\Gin@nat@height\fi}
\makeatother
\graphicspath{{./}{vignettes/manuscript/}{assets/}}
\DeclareRobustCommand{\pkg}[1]{\textsf{#1}}
\DeclareRobustCommand{\code}[1]{\nolinkurl{#1}}
\setlength{\emergencystretch}{2em}
\setkeys{Gin}{width=\linewidth,height=0.82\textheight,keepaspectratio}

%% ---- JSS metadata ----
\author{Arnab Aich\\Florida State University}
\title{CausalMixGPD: Bayesian NP Model with Optional Heavy-Tail Augmentaion and Causal Extensions}
\Plainauthor{Arnab Aich}
\Plaintitle{CausalMixGPD: Bayesian Conditional Dirichlet Process Mixtures with Generalized Pareto Tails and Causal Extensions}
\Shorttitle{CausalMixGPD}

\Abstract{
Heavy-tailed outcomes arise in many scientific domains where rare events dominate risk, cost, or policy decisions. Classical parametric models often fit the center of the distribution at the expense of tail misspecification, while purely nonparametric methods typically lack principled extrapolation in regions with sparse data. Extreme value theory (EVT) provides asymptotic justification for generalized Pareto distributions (GPDs) for threshold exceedances, whereas Bayesian nonparametric Dirichlet process mixtures (DPMs) provide flexible modeling for heterogeneous bulk behavior. This article introduces \texttt{CausalMixGPD}, an \textsf{R} package implementing a Bayesian semiparametric framework that combines conditional DPMs for the bulk distribution with covariate-dependent GPD tails for exceedances, with posterior predictive inference for densities, distribution functions, quantiles, survival probabilities, and tail risks. The framework is conditional by construction; unconditional models arise as a special case when covariates are omitted. The package supports both Chinese restaurant process (CRP) and stick-breaking (SB) representations of the Dirichlet process, multiple kernel families for bulk modeling, and an extension to causal inference via treatment-specific conditional distributions that enable inference on distributional and tail-sensitive causal estimands.
}

\Keywords{Bayesian nonparametrics; Dirichlet Process Mixture; Generalized Pareto Distribution; Extreme Value Theory; Causal Inference; \textsf{R}; \texttt{nimble}}
\Plainkeywords{Bayesian nonparametrics, Dirichlet process mixture, generalized Pareto distribution, extreme value theory, posterior prediction, causal inference, quantile treatment effects, R, nimble}

%% ---- Begin document ----
\IfFileExists{upquote.sty}{\usepackage{upquote}}{}
\begin{document}
\maketitle
\sloppy



\section{Introduction}
\label{sec:intro}

Causal inference is often summarized through average effects, but many
applications require distributional comparisons between treatment
regimes. In clinical trials, investigators may want to understand how a
treatment changes outcomes across the outcome scale, including effects
among individuals in the upper tail where events are rare but
consequences can be substantial. This shift from mean-based to
distribution-based causal questions makes causal analysis tightly linked
to how well we can estimate (and predict) outcome distributions under
competing treatment assignments.

A growing \textsf{R} ecosystem supports Bayesian causal inference
through flexible outcome modeling and posterior uncertainty
quantification. For example,
\href{https://cran.r-project.org/package=bcf}{\textsf{bcf}} implements
Bayesian Causal Forests for binary treatments and continuous outcomes,
providing posterior inference for both average and heterogeneous
treatment effects \citep{bcf_pkg}. The package
\href{https://cran.r-project.org/package=bartCause}{\textsf{bartCause}}
builds causal estimators around Bayesian additive regression trees,
using BART as a regression engine for treatment-response modeling and
causal effect estimation \citep{bartCause_pkg,hill2011bart}.
Complementing tree-based approaches,
\href{https://cran.r-project.org/package=BCEE}{\textsf{BCEE}} implements
a Bayesian model-averaging strategy for causal effect estimation with
support for binary or continuous exposures and outcomes
\citep{BCEE_pkg,talbot2015bcee}. A recurring practical difficulty in
quantile-based causal reporting is quantile crossing issue, i.e., when
two quantile curves for different indices cross each other. To avoid
that issue this package (CausalMixGPD) directly model outcome densities and derives
quantiles and other functionals from the posterior predictive distributions.

In the context of density estimation and conditional density estimation, DPMs provide a
versatile non parametric solution. It can identify multimodality and latent heterogeneity without specifying number of latent
mixing components.
 In expreimental framework \textsf{R} is
well supported by Bayesian nonparametric mixture modeling software based
on DP ideas. For example,
\href{https://cran.r-project.org/package=dirichletprocess}{\textsf{dirichletprocess}}
provides a flexible interface for specifying DPMs and
fitting them by posterior simulation, yielding posterior inference for
mixture weights, component allocations, and predictive densities under
user-chosen kernels \citep{dirichletprocess_pkg}. The package
\href{https://cran.r-project.org/package=DPpackage}{\textsf{DPpackage}}
offers a broad suite of Bayesian semi- and nonparametric models in
\textsf{R}, including DP-based formulations for marginal and conditional
densities, regression functions, and model-based predictive
distributions \citep{jara2011DPpackage}. More recently,
\href{https://cran.r-project.org/package=BNPmix}{\textsf{BNPmix}}
emphasizes practical Bayesian nonparametric mixture approaches,
providing posterior inference for clustering structure and
density/regression estimation under DP and Pitman--Yor-type mixture
specifications \citep{canale2021BNPmix}. However, DPMs are
suitable for modeling the bulk of the distribution but do not explicitly
model tail behavior, which can lead to poor estimation and prediction in
the tails when data are sparse or a heavy tail is present. mMoreover, these packages
use a single kernel family for the entire distribution,
 which may not be ideal when the bulk and tail exhibit different characteristics.

Bayesian implementations of EVT are available in
\textsf{R} when posterior uncertainty quantification for tail
quantitites are required. For example,
\href{https://cran.r-project.org/package=evdbayes}{\textsf{evdbayes}}
provides Bayesian inference for standard extreme-value models via
posterior simulation, enabling posterior and posterior-predictive
summaries for extreme-event quantities \citep{evdbayes_pkg}. For
peaks-over-threshold (POT) approaches, the package
\href{https://cran.r-project.org/package=POT}{\textsf{POT}} implements a
focused collection of tools for GPD modeling and threshold-based tail
analysis \citep{POT_pkg}. Finally, for extreme quantile estimation and
rare-event tail summaries,
\href{https://cran.r-project.org/package=extremefit}{\textsf{extremefit}}
provides a published software implementation targeted specifically at
extreme-quantile estimation \citep{durrieu2018extremefit}.

This article introduces, an \textsf{R} package (\pkg{CausalMixGPD}) that
connects Bayesian nonparametric bulk modeling with extreme-value tail
modeling, while also supporting robust causal analyses in presence of covariates(confounders).
The package provides two complementary capabilities:

\begin{itemize}
   \item \textbf{Bayesian Non-Parametric density modeling and prediction with optional GPD tail:}
   \pkg{CausalMixGPD} fits DPMs using a user-chosen kernel below a certain threshold (bulk part),
   and can optionally splice a generalized Pareto tail above the threshold (tail part) to enable the estimation of tail quantities
   and posterior predictive inference for rare-event summaries.
   \item \textbf{Causal estimation and prediction for RCTs and observational studies:}
   This package also supports causal analyses by modeling treatment-specific potetntial outcome
   densities and reporting contrasts as functionals of the posterior over densities.
 \end{itemize}

Taken together, \pkg{CausalMixGPD} offers a coherent workflow for
density estimation, prediction, and distributional causal reporting when
tail behavior is scientifically central. The package is implemented in \textsf{R} with computationally intensive components
coded in \textsf{C++} and integrated via the \texttt{nimble} framework for flexible model specification and
efficient posterior simulation \citep{nimble_pkg}.
The remainder of this article is organized as follows: Section~\ref{sec:prelim} provides background on the key
modeling components, including DPMs, GPDs,
and the spliced bulk-tail construction.
Section~\ref{sec:overview} gives an overview of the package's main functions and user interface.
Section~\ref{sec:app-causal} and Section~\ref{sec:app-clustering} illustrates the package's capabilities through applications to causal inference and clustering, respectively.
 Finally, Section~\ref{sec:discussion} concludes with a discussion of limitations and future directions.

\section{Background and Preliminaries}
\label{sec:prelim}

This package treats conditional outcome distributions as the fundamental
inferential object. All summaries reported later are functionals of
estimated conditional distributions, including conditional quantiles,
tail probabilities, and distributional causal contrasts.

Let the outcome--covariate pair be defined by \begin{equation}
(Y,X) \in \mathbb{R}\times \mathbb{R}^p .
\end{equation} For a fixed covariate value $x$, the conditional
cumulative distribution function (CDF) and density are \begin{equation}
F(y\mid x)=\Pr\;(Y\le y\mid X=x),
\qquad
f(y\mid x)=\frac{\partial}{\partial y}F(y\mid x),
\end{equation} and the conditional quantile function is defined by
\begin{equation}
Q(\tau\mid x)=\inf\Bigl\{y\in\mathbb{R}:\,F(y\mid x)\ge \tau\Bigr\},
\qquad 0<\tau<1.
\end{equation}



\subsection{DPM Model for Conditional Densities}
\label{subsec:dpm}

I model the bulk ($Y<u$) conditional density of $Y$ given $X=x$ using a
DPM regression model. Let $k(\cdot\mid x;\theta)$
be a user-chosen kernel family indexed by parameter $\theta$. The
hierarchical specification is \begin{equation}
Y \mid X=x,\theta \sim k(\cdot\mid x;\theta),
\qquad
\theta \mid H \sim H,
\qquad
H \mid \kappa,H_0 \sim \mathrm{DP}(\kappa,H_0),
\label{eq:dpm_hier}
\end{equation} where $H$ is an unknown mixing distribution on the
kernel-parameter space, $H_0$ is a base measure, and $\kappa>0$ is a
concentration parameter \citep{Ferguson1973,Lo1984}.

Marginalizing $H$ yields an (almost surely) discrete mixing distribution
and an infinite mixture representation: \begin{equation}
f_{\mathrm{DP}}(y\mid x)
=
\int k(y\mid x;\theta)\,dH(\theta)
=
\sum_{j=1}^{\infty} w_j\,k(y\mid x;\theta_j),
\qquad
w_j\ge 0,
\qquad
\sum_{j=1}^{\infty} w_j=1.
\label{eq:dpm_mixture}
\end{equation}

A constructive representation for the weights is given by Sethuraman's
stick-breaking construction \citep{Sethuraman1994}: \begin{equation}
w_1=V_1,
\qquad
w_j=V_j\prod_{\ell<j}(1-V_\ell),
\qquad
V_j \stackrel{\mathrm{iid}}{\sim}\mathrm{Beta}(1,\kappa),
\qquad
\theta_j \stackrel{\mathrm{iid}}{\sim} H_0.
\label{eq:sb}
\end{equation} An alternative marginal representation is obtained by
integrating out $H$, which induces the Blackwell--MacQueen P'olya urn
predictive scheme (often described via the Chinese restaurant process)
\citep{BlackwellMacQueen1973,Neal2000}: \begin{equation}
\theta_{\text{new}} \mid \theta_{1:n}
\sim
\frac{\kappa}{\kappa+n}H_0(\cdot)
+
\frac{1}{\kappa+n}\sum_{m=1}^{n}\delta_{\theta_m}(\cdot).
\label{eq:crp}
\end{equation} Both stick-breaking (blocked/truncated) and
allocation-based representations are widely used for posterior
simulation in DPM models \citep{EscobarWest1995,Neal2000}.  DPM modeling is flexible, it works without specifying
the number of mixture components, and can capture complex distributional features such as
multimodality, skewness, and heteroscedasticity by choosing an appropriate kernel family.
This makes DPMs a natural choice for modeling the bulk of the conditional distribution, where data are abundant
and complex features may be present. The choice of kernel family is user-controlled
which I will discuss in Section~\ref{sec:overview}.


\subsection{GPD for Threshold Exceedances}
\label{subsec:gpd}

At the tail region EVT provides a canonical approximation for exceedances
above a high threshold. Let $u(x)$ be a (possibly covariate-dependent)
high threshold and define the exceedance \begin{equation}
Z = Y-u(x),
\end{equation} considered conditionally on the event $Y>u(x)$. The
conditional excess distribution function is \begin{equation}
F_u(z\mid x)=\Pr\;\!\bigl(Y-u(x) \le z \mid Y>u(x),X=x\bigr),
\qquad z\ge 0.
\label{eq:cedf}
\end{equation} Under broad regularity conditions, $F_u(\cdot)$ is
well-approximated by a GPD for
sufficiently large $u$ \citep{BalkemaDeHaan1974,Pickands1975,Coles2001}.
With scale parameter $\sigma(x)>0$ and shape parameter
$\xi\in\mathbb{R}$, the CDF for tail excedances is given by
\begin{equation}
F_{\mathrm{GPD}}(y\mid u(x),\sigma(x))=
\begin{cases}
1-\left(1+\xi\dfrac{y-u(x)}{\sigma(x)}\right)^{-1/\xi}, & \xi\neq 0,\\[8pt]
1-\exp\!\left(-\dfrac{y-u(x)}{\sigma(x)}\right), & \xi=0,
\end{cases}
\qquad y>u(x),
\label{eq:gpd_cdf}
\end{equation} with support $y\ge u(x)$ when $\xi\ge 0$, and
$u(x)\le y\le u(x)-\sigma(x)/\xi$ when $\xi<0$. The corresponding
density is \begin{equation}
f_{\mathrm{GPD}}(y\mid u(x),\sigma(x),\xi)=
\begin{cases}
\dfrac{1}{\sigma(x)}\left(1+\xi\dfrac{y-u(x)}{\sigma(x)}\right)^{-1/\xi-1}, & \xi\neq 0,\\[10pt]
\dfrac{1}{\sigma(x)}\exp\!\left(-\dfrac{y-u(x)}{\sigma(x)}\right), & \xi=0.
\end{cases}
\label{eq:gpd_pdf}
\end{equation} The associated GPD quantile function is given by
 \begin{equation}
Q_{\mathrm{GPD}}(\tau\mid u(x),\sigma(x),\xi)=
\begin{cases}
u(x)-\sigma(x)\log(1-\tau), & \xi=0,\\[6pt]
u(x)+\dfrac{\sigma(x)}{\xi}\Bigl((1-\tau)^{-\xi}-1\Bigr), & \xi\neq 0,
\end{cases}
\qquad 0<\tau<1.
\label{eq:gpd_quantile}
\end{equation}



\begin{knitrout}\footnotesize
\definecolor{shadecolor}{rgb}{0.969, 0.969, 0.969}\color{fgcolor}

{\centering \includegraphics[width=0.82\linewidth]{cache/figs/gpdDensityFigure} 

}


\end{knitrout}

The figure above illustrates the GPD density for different shape parameters $\xi$ with a fixed
scale parameter $\sigma$ and threshold $u$. The shape parameter $\xi$ controls the tail heaviness: $\xi < 0$
corresponds to a light tail with a finite upper bound, $\xi = 0$ corresponds to an exponential tail, and $\xi > 0$
corresponds to a heavy tail with infinite support. The GPD provides a flexible family for modeling exceedances
above a high threshold, making it a natural choice for tail modeling in the spliced construction described next.


\subsection{Spliced DPM--GPD Conditional Distribution}
\label{subsec:splice}

DPM regression models~Eq:~\eqref{eq:dpm_hier} provide a
flexible representation of the \emph{bulk} of the conditional
distribution of $Y\mid X=x$, however, in the far
upper tail-where data are intrinsically sparse, EVT provides an parametric description of
threshold exceedances via the GPD in Eq:~\eqref{eq:gpd_pdf}. I therefore adopt a
\emph{spliced} (bulk--tail) construction that retains the DPM fit below
a high threshold and replaces the upper tail by a GPD component
\citep{Frigessi2002,Behrens2004,Coles2001}. Let \begin{equation}
p_u(x)=F_{\mathrm{DP}}(u(x)\mid x;\Theta)
\label{eq:pu_theta}
\end{equation} denote the bulk probability mass below a threshold ($u(x)$)
under the DPM fit. The spliced conditional CDF is given as \begin{equation}
F(y\mid x;\Theta,\Phi)=
\begin{cases}
F_{\mathrm{DP}}(y\mid x;\Theta), & y \le u(x),\\[6pt]
p_u(x;\Theta)
+
\bigl\{1-p_u(x;\Theta)\bigr\}\,
F_{\mathrm{GPD}}\!\bigl(y\mid x;\Phi\bigr),
& y>u(x),
\end{cases}
\label{eq:splice_cdf}
\end{equation} where  $F_{\mathrm{DP}}(\cdot)$ is the DPM CDF on the bulk region and
 $F_{\mathrm{GPD}}(\cdot)$ is the GPD CDF on the
exceedance region (Section~\ref{subsec:gpd}). The CDF~\eqref{eq:splice_cdf} is
continuous at $u(x)$ and yields a proper distribution. Differentiating \eqref{eq:splice_cdf} yields the spliced density
\begin{equation}
f(y\mid x;\Theta,\Phi)=
\begin{cases}
f_{\mathrm{DP}}(y\mid x;\Theta), & y \le u(x),\\[8pt]
\bigl\{1-p_u(x;\Theta)\bigr\}\,
f_{\mathrm{GPD}}\!\bigl(y\mid x;\Phi\bigr),
& y>u(x),
\end{cases}
\label{eq:splice_pdf}
\end{equation} so the tail density is a \emph{scaled} GPD density, with
scaling ensuring normalization of $f(\cdot\mid x;\Theta,\Phi)$.

For $y>u(x)$, the conditional survival function decomposes as
\begin{equation}
\Pr\;(Y>y\mid X=x;\Theta,\Phi)
=
\bigl\{1-p_u(x;\Theta)\bigr\}
\Bigl[1-F_{\mathrm{GPD}}\!\bigl(y\mid x;\Phi\bigr)\Bigr].
\label{eq:splice_tailprob}
\end{equation} Quantiles follow by inversion of \eqref{eq:splice_cdf}.
For $\tau\in(0,1)$, \begin{equation}
Q(\tau\mid x;\Theta,\Phi)=
\begin{cases}
Q_{\mathrm{DP}}(\tau\mid x;\Theta), & \tau \le p_u(x;\Theta),\\[8pt]
Q_{\mathrm{GPD}}\!\bigl(\tilde{\tau}(x;\Theta)\mid x;\Phi\bigr), & \tau>p_u(x;\Theta),
\end{cases}
\label{eq:splice_quantile_case}
\end{equation} where the tail-rescaled probability level is
\begin{equation}
\tilde{\tau}(x;\Theta)=\dfrac{\tau-p_u(x;\Theta)}{1-p_u(x;\Theta)}\in(0,1),
\label{eq:tau_tilde}
\end{equation} and $Q_{\mathrm{GPD}}$ is given in
\eqref{eq:gpd_quantile}. Thus, upper quantiles are determined jointly by
the bulk exceedance probability $1-p_u(x;\Theta)$ and the conditional
exceedance distribution governed by $\Phi$.



\begin{knitrout}\footnotesize
\definecolor{shadecolor}{rgb}{0.969, 0.969, 0.969}\color{fgcolor}

{\centering \includegraphics[width=0.82\linewidth]{cache/figs/splicedDensityFigure} 

}


\end{knitrout}


To make the splicing and downstream functionals explicit, we collect the
unknown bulk regression parameters into $\Theta$ and the tail parameters
into $\Phi$. Specifically, the bulk model in \eqref{eq:dpm_mixture} can
be written as \begin{equation}
f_{\mathrm{DP}}(y\mid x)
=
\sum_{j=1}^{\infty} w_j \, k\!\left(y\mid x;\theta_j(x)\right),
\qquad
\Theta=\{\beta_j\}_{j\ge 1},
\label{eq:bulk_theta}
\end{equation} where $\beta_j$ denotes component-specific
regression coefficients and $$\theta_j(x) = g(x^{\top} \beta_j)$$ denotes the kernel
parameter vector induced at covariate value $x$ through the
corresponding coefficients in $\Theta$ by a user chosen link function $g(\cdot)$. Mixture weights $\{w_j\}$ are
generated by the DP backend (SB or CRP representation) which we will discuss in Section~\eqref{sec:overview}.
It is to be noted that $\Theta$ may contain other kernel
parameters that may not depend on covariates. As an example if the kernel is a Gaussian density, then $\Theta$ may
 contain regression coefficients for the mean ($\beta^{\mu}_j$) and a variance parameter ($\sigma_j$) that does not depend on covariates.


For the tail block, let \begin{equation}
\Phi=\{\beta_u,\beta_\sigma,\xi\},
\label{eq:phi_def}
\end{equation} where $u(x)$ and $\sigma(x)$ are induced by covariates
and tail-regression coefficients as
$$u(x) = g_u(x^{\top} \beta_u), \qquad \sigma(x) = g_\sigma(x^{\top} \beta_\sigma),$$
 with $g_u(\cdot)$ and $g_\sigma(\cdot)$ being user-chosen link functions.
Lastly, $\xi$ is the tail shape parameter as discussed in Section~\eqref{subsec:gpd}.
Thus, $\Theta$ and $\Phi$ are global parameter collections, while $x$
enters through these mappings. The package allows users to control which
kernel and tail parameters vary with covariates (e.g., via link-based
regression for selected parameters) and which are treated as fixed or
purely a priori parameters (Section~\ref{sec:overview}).


This spliced conditional distribution is the central predictive object
used throughout the package. In the causal extension
(Section~\ref{subsec:causal}), the same construction is applied
separately to treatment-specific conditional distributions
$F_a(\cdot\mid x)$, enabling causal contrasts based on densities,
quantiles, and tail risks rather than only mean-scale effects.

\subsection{Causal Inference: Notation, Identification, and Reported Estimands}
\label{subsec:causal}

I use the potential outcomes framework
\citep{Neyman1923,Rubin1974,ImbensRubin2015}. Let $A$ denote a binary
treatment indicator, let $X$ denote observed pre-treatment covariates,
and let $Y$ denote the observed outcome. Each unit has two potential
outcomes, $Y(1)$ and $Y(0)$, and the observed outcome satisfies
\begin{equation}
Y = A\,Y(1) + \bigl(1-A\bigr)\,Y(0).
\label{eq:po_observed}
\end{equation}

The treatment-specific conditional distribution functions are the
fundamental inferential objects: \begin{equation}
F_a(y\mid x)=\Pr\;\{Y(a)\le y \mid X=x\},\qquad a\in\{0,1\}.
\label{eq:po_cond_cdf}
\end{equation} In this package, the corresponding conditional densities
are \begin{equation}
f_a(y\mid x)=\frac{\partial}{\partial y}F_a(y\mid x)
\label{eq:po_cond_pdf}
\end{equation} where $f_a(y \mid x)$ assumes the spliced DPM structure
in~\eqref{eq:splice_pdf}. Conditional quantiles are defined by
inversion of the conditional CDF: \begin{equation}
Q_a(\tau\mid x)=\inf\{y:\,F_a(y\mid x)\ge \tau\},\qquad \tau\in(0,1).
\label{eq:po_cond_quantile}
\end{equation} Conditional means are functionals of the conditional
density: \begin{equation}
\mu_a(x)=\mathrm{E}\{Y(a)\mid X=x\}=\int y\, f_a(y\mid x)\,dy.
\label{eq:po_cond_mean}
\end{equation}

Marginal (standardized) distributions are obtained by averaging
conditional distributions over a covariate distribution. The package
reports population-standardized and treated-standardized functionals:
\begin{equation}
F_a^{\mathrm{m}}(y)=\mathrm{E}\bigl\{F_a(y\mid X)\bigr\},
\qquad
F_a^{\mathrm{t}}(y)=\mathrm{E}\bigl\{F_a(y\mid X)\mid A=1\bigr\}.
\label{eq:marginal_cdf_defs}
\end{equation} The corresponding marginal quantiles are \begin{equation}
Q_a^{\mathrm{m}}(\tau)=\inf\{y:\,F_a^{\mathrm{m}}(y)\ge \tau\},
\qquad
Q_a^{\mathrm{t}}(\tau)=\inf\{y:\,F_a^{\mathrm{t}}(y)\ge \tau\}.
\label{eq:marginal_quantile_defs}
\end{equation} Marginal means are \begin{equation}
\mu_a^{\mathrm{m}}=\mathrm{E}\bigl\{\mu_a(X)\bigr\},
\qquad
\mu_a^{\mathrm{t}}=\mathrm{E}\bigl\{\mu_a(X)\mid A=1\bigr\}.
\label{eq:marginal_mean_defs}
\end{equation}

Identification of these quantities depends on the study design. In
randomized experiments, randomization together with positivity implies
exchangeability between treatment assignment and potential outcomes
\citep{ImbensRubin2015}. In observational studies, we rely on the
following conditions \citep{RosenbaumRubin1983,ImbensRubin2015}:

\begin{itemize}
\item \textbf{SUTVA:} no interference between units and no hidden versions of treatment.
\item \textbf{Consistency:} if $A=a$, then the observed outcome equals the corresponding potential outcome, $Y=Y(a)$.
\item \textbf{Conditional ignorability:} treatment is as-if randomized given covariates,
\begin{equation}
\bigl(Y(0),Y(1)\bigr)\ \perp\ A \ \mid\ X .
\label{eq:ignorability}
\end{equation}
\item \textbf{Overlap (positivity):} each covariate profile has positive probability of receiving either treatment,
\begin{equation}
0<\Pr\;(A=1\mid X=x)<1
\label{eq:overlap}
\end{equation}
for covariate values of interest.
\end{itemize}

Under these conditions, the treatment-specific conditional distributions
are identified from the observed data: \begin{equation}
F_a(y\mid x)=\Pr\;(Y\le y\mid A=a,X=x),\qquad a\in\{0,1\}.
\label{eq:identification_cdf}
\end{equation}

In observational settings, we use the propensity score (PS) to support
the design stage \citep{RosenbaumRubin1983}. The PS is the treatment
assignment probability given covariates: \begin{equation}
\rho(x)=\Pr\;(A=1\mid X=x).
\label{eq:ps_def}
\end{equation} A key property of the true propensity score is that
conditioning on it balances covariates across treatment groups:
\begin{equation}
A\ \perp\ X\ \mid\ \rho(X).
\label{eq:ps_balance}
\end{equation} With estimated scores this balancing property is
approximate in finite samples. Following the two-stage strategy
described in the thesis, we first estimate $\rho(x)$ using only $(A,X)$
and then augment the outcome model covariates with the estimated score
\citep{RosenbaumRubin1983}: \begin{equation}
\widehat{\rho}(x)\approx \rho(x),
\qquad
r(x)=\bigl(x^{\mathsf{T}},\widehat{\rho}(x)\bigr)^{\mathsf{T}}.
\label{eq:augmented_covariate}
\end{equation} The conditional outcome distributions are then modeled
separately within each treatment arm as functions of the augmented
predictor $r(x)$, and causal contrasts are computed as functionals of
the estimated treatment-specific distributions.

The package reports the following causal estimands:

\begin{itemize}
\item \textbf{Average Treatment Effect (ATE):}
\begin{equation}
\mathrm{ATE}=\mu_1^{\mathrm{m}}-\mu_0^{\mathrm{m}}.
\label{eq:ate}
\end{equation}
\item \textbf{Average Treatment Effect on the Treated (ATT):}
\begin{equation}
\mathrm{ATT}=\mu_1^{\mathrm{t}}-\mu_0^{\mathrm{t}}.
\label{eq:att}
\end{equation}
\item \textbf{Quantile Treatment Effect (QTE):} for any $\tau\in(0,1)$ \citep{Firpo2007},
\begin{equation}
\mathrm{QTE}(\tau)=Q_1^{\mathrm{m}}(\tau)-Q_0^{\mathrm{m}}(\tau).
\label{eq:qte}
\end{equation}
\item \textbf{Quantile Treatment Effect on the Treated (QTT):} for any $\tau\in(0,1)$,
\begin{equation}
\mathrm{QTT}(\tau)=Q_1^{\mathrm{t}}(\tau)-Q_0^{\mathrm{t}}(\tau).
\label{eq:qtt}
\end{equation}
\item \textbf{Conditional Average Treatment Effect (CATE):}
\begin{equation}
\mathrm{CATE}(x)=\mu_1(x)-\mu_0(x).
\label{eq:cate}
\end{equation}
\item \textbf{Conditional Quantile Treatment Effect (CQTE):} for any $\tau\in(0,1)$,
\begin{equation}
\mathrm{CQTE}(\tau\mid x)=Q_1(\tau\mid x)-Q_0(\tau\mid x).
\label{eq:cqte}
\end{equation}
\end{itemize}

This section established the main probabilistic and causal building
blocks used in the remainder of the work. I first defined conditional
distributional targets and introduced a DPM model
to flexibly estimate bulk conditional outcome densities. I then
reviewed generalized Pareto modeling for threshold exceedances and
constructed a spliced DPM--GPD conditional distribution that preserves
the estimated bulk while enabling principled tail extrapolation.
Finally, we introduced the potential outcomes framework and the
treatment-effect estimands reported in the package, showing how marginal
and conditional effects can be obtained as functionals of
treatment-specific conditional distributions (with propensity-score
augmentation supporting confounding adjustment in observational
studies). Subsequent chapters build on these preliminaries to specify
the full hierarchical model, posterior computation, and empirical
evaluation of tail-aware distributional causal effects.

\section{Package Overview: Workflow, Model, and Inference}
\label{sec:overview}

This section describes the main workflow provided by \pkg{CausalMixGPD}
for one-arm conditional modeling and prediction when a generalized
Pareto tail is spliced above a covariate-dependent threshold. The causal
extension reuses the same modeling, estimation, and prediction machinery
after introducing treatment-specific outcome models and a propensity
score component; we defer those additions to the causal section.

\subsection{Augmented DPM--GPD model, likelihood, and priors}
\label{subsec:overview-model}

Let the observed data be $$
  \mathcal{D}=\{(y_i,x_i)\}_{i=1}^n,
$$

\begin{knitrout}\footnotesize
\definecolor{shadecolor}{rgb}{0.969, 0.969, 0.969}\color{fgcolor}\begin{kframe}
\begin{alltt}
\hlkwd{data}\hldef{(}\hlsng{"nc_posX100_p3_k2"}\hldef{,} \hlkwc{package} \hldef{=} \hlsng{"CausalMixGPD"}\hldef{)}
\hldef{dat} \hlkwb{<-} \hlkwd{data.frame}\hldef{(}\hlkwc{y} \hldef{= nc_posX100_p3_k2}\hlopt{$}\hldef{y, nc_posX100_p3_k2}\hlopt{$}\hldef{X)}
\end{alltt}
\end{kframe}
\end{knitrout}

with $x_i\in\mathbb{R}^p$ denoting the covariate vector (including an
intercept if desired). The target conditional distribution
$F(\cdot\mid x)$ and its spliced DPM--GPD construction are defined in
Section~\ref{sec:prelim}, with the spliced CDF and density given in
\eqref{eq:splice_cdf}--\eqref{eq:splice_pdf}. Posterior computation in
\pkg{CausalMixGPD} approximates the bulk DPM by a finite truncation with
$J$ mixture components, $$
  f_{\mathrm{DP}}(y\mid x;\Theta)\;\approx\;\sum_{j=1}^{J} w_j\,k(y\mid x;\theta_j),
  \qquad j=1,\ldots,J,
$$ where $j$ indexes mixture components, $w_j\ge 0$, and
$\sum_{j=1}^J w_j=1$. Kernel parameters can be constant or
covariate-dependent through link specifications of the form
\begin{equation}
  g\!\left(\psi_j(x)\right)=x^\top \beta_{j,\psi},
  \label{eq:overview-link}
\end{equation} where $g(\cdot)$ enforces parameter support (e.g.,
positive scale parameters).

\begin{knitrout}\footnotesize
\definecolor{shadecolor}{rgb}{0.969, 0.969, 0.969}\color{fgcolor}\begin{kframe}
\begin{alltt}
\hldef{mcmc_fixed} \hlkwb{<-} \hlkwd{list}\hldef{(}
  \hlkwc{niter} \hldef{=} \hlnum{200}\hldef{,}
  \hlkwc{nburnin} \hldef{=} \hlnum{50}\hldef{,}
  \hlkwc{thin} \hldef{=} \hlnum{1}\hldef{,}
  \hlkwc{nchains} \hldef{=} \hlnum{1}\hldef{,}
  \hlkwc{seed} \hldef{=} \hlnum{1}
\hldef{)}
\end{alltt}
\end{kframe}
\end{knitrout}

\begin{knitrout}\footnotesize
\definecolor{shadecolor}{rgb}{0.969, 0.969, 0.969}\color{fgcolor}\begin{kframe}
\begin{alltt}
\hldef{fit} \hlkwb{<-} \hlkwd{dpmgpd}\hldef{(}
  \hlkwc{formula} \hldef{= y} \hlopt{~} \hldef{x1} \hlopt{+} \hldef{x2} \hlopt{+} \hldef{x3,}
  \hlkwc{data} \hldef{= dat,}
  \hlkwc{backend} \hldef{=} \hlsng{"sb"}\hldef{,}
  \hlkwc{kernel} \hldef{=} \hlsng{"lognormal"}\hldef{,}
  \hlkwc{components} \hldef{=} \hlnum{5}\hldef{,}
  \hlkwc{mcmc} \hldef{= mcmc_fixed}
\hldef{)}
\end{alltt}
\end{kframe}
\end{knitrout}



\begin{knitrout}\footnotesize
\definecolor{shadecolor}{rgb}{0.969, 0.969, 0.969}\color{fgcolor}\begin{kframe}
\begin{verbatim}
MixGPD fit | backend: Stick-Breaking Process | kernel:
Lognormal Distribution | GPD tail: TRUE
n = 100 | components = 5 | epsilon = 0.025
MCMC: niter=200, nburnin=50, thin=1, nchains=1
Fit
Use summary() for posterior summaries; plot() for
diagnostics; predict() for predictions. 
\end{verbatim}
\end{kframe}
\end{knitrout}

For the spliced model, let $u(x)$ be the threshold, $\sigma(x)>0$ the
tail scale, and $\xi\in\mathbb{R}$ the tail shape (as in
Section~\ref{subsec:splice}). Define
$\delta_i=\mathbf{1}\{y_i>u(x_i)\}$. Using the spliced density
\eqref{eq:splice_pdf} with the truncated bulk mixture, the likelihood is
\begin{equation}
  L(\Theta,\Phi;\mathcal{D})
  =\prod_{i=1}^n f(y_i\mid x_i;\Theta,\Phi),
  \label{eq:overview-lik}
\end{equation} where $\Theta$ collects the truncation-level bulk
parameters $\{w_j,\theta_j\}_{j=1}^J$, and $\Phi$ collects the tail
parameters governing $u(x)$, $\sigma(x)$, and $\xi$. Equivalently, the
log-likelihood can be written as \begin{align}
  \log L(\Theta,\Phi;\mathcal{D})
  \;=&\;\sum_{i=1}^n (1-\delta_i)\,\log f_{\mathrm{DP}}(y_i\mid x_i;\Theta)
  \nonumber\\
  &\;+\;\sum_{i=1}^n \delta_i
  \Bigl[
    \log\!\bigl\{1-F_{\mathrm{DP}}(u(x_i)\mid x_i;\Theta)\bigr\}
    +\log f_{\mathrm{GPD}}\!\bigl(y_i\mid x_i;\Phi\bigr)
  \Bigr],
  \label{eq:overview-loglik}
\end{align} where $f_{\mathrm{GPD}}$ is given in \eqref{eq:gpd_pdf} and
$F_{\mathrm{DP}}(\cdot\mid x)$ is the bulk mixture CDF induced by the
truncated mixture.

Priors follow the decomposition implied by \eqref{eq:sb} and
\eqref{eq:crp}. Under the SB backend, the
truncation-level weights are generated through Beta stick breaks
(Section~\ref{sec:prelim}, \eqref{eq:sb}), which are represented in
\pkg{nimble} using the Beta distribution in model code (see
\code{dbeta}). Under the CRP backend, the
allocation-based representation is implemented using the CRP
distribution in \pkg{nimble} (see \code{dCRP}), matching the predictive
representation in \eqref{eq:crp}. Kernel parameters and tail parameters
are assigned weakly informative priors, with regression coefficients
$\beta_{j,\psi}$ arising through the link representation
\eqref{eq:overview-link}. The same link construction is used for
covariate-dependent $u(x)$ and $\sigma(x)$ when those are modeled as
functions of $x$.

\subsection{MCMC algorithm, convergence diagnostics, and parameter extraction}
\label{subsec:overview-mcmc}

Posterior simulation targets
$p(\Theta,\Phi\mid \mathcal{D}) \propto L(\Theta,\Phi;\mathcal{D})\,p(\Theta,\Phi)$
using the likelihood \eqref{eq:overview-lik} and the prior structure
implied by the SB or CRP representation (Section~\ref{sec:prelim},
\eqref{eq:sb}--\eqref{eq:crp}). Given a model specification,
\pkg{CausalMixGPD} runs MCMC through \pkg{nimble}
\citep{devalpine2017nimble,nimble_pkg} under the chosen backend; SB uses
Beta stick breaks in model code (see \code{dbeta}) and CRP uses the CRP
distribution for allocations (see \code{dCRP}). Users control the number
of iterations, burn-in, thinning, and number of chains through the
\code{mcmc} argument in the wrappers.

Convergence checking is supported through standard multi-chain
diagnostics and trace-based visualizations. The model object supports
\code{summary()} for posterior summaries, \code{plot()} for diagnostic
plots (trace, running mean, autocorrelation, and multi-chain diagnostics
such as $\widehat{R}$ when multiple chains are used), and
\code{params()} as a lightweight extractor returning posterior mean
parameters reshaped to natural dimensions. Interval summaries used by
\code{predict()} can be chosen as equal-tailed credible intervals or HPD
intervals as described in Section~\ref{subsec:overview-estimation}
\citep{Gelman2013BDA,coda_pkg}. Diagnostic plotting relies on
established MCMC diagnostic tooling \citep{ggmcmc_pkg}.

\begin{knitrout}\footnotesize
\definecolor{shadecolor}{rgb}{0.969, 0.969, 0.969}\color{fgcolor}\begin{kframe}
\begin{alltt}
\hlkwd{print}\hldef{(fit)}
\hldef{p_hat} \hlkwb{<-} \hlkwd{params}\hldef{(fit)}
\hldef{p_hat}
\end{alltt}
\end{kframe}
\end{knitrout}



\begin{knitrout}\footnotesize
\definecolor{shadecolor}{rgb}{0.969, 0.969, 0.969}\color{fgcolor}\begin{kframe}
\begin{verbatim}
MixGPD fit | backend: Stick-Breaking Process | kernel:
Lognormal Distribution | GPD tail: TRUE
n = 100 | components = 5 | epsilon = 0.025
MCMC: niter=200, nburnin=50, thin=1, nchains=1
Fit
Use summary() for posterior summaries; plot() for
diagnostics; predict() for predictions. 
Posterior mean parameters

$alpha
[1] 2.034

$w
[1] 0.4224 0.3063 0.1465 0.1158

$beta_meanlog
            x1      x2       x3
comp1  0.34970 -0.5504 -0.31760
comp2 -0.09937 -0.9372  0.10770
comp3  0.22010  0.5477  0.70350
comp4  0.04718  0.3581 -0.30940
comp5 -0.38810  0.8843 -0.04477

$sdlog
[1] 1.517 2.222 1.481 1.482

$beta_threshold
[1] -0.05201 -0.10530  0.07445

$sdlog_u
[1] 1.356

$beta_tail_scale
[1] 0.28030 0.04619 0.07101

$tail_shape
[1] 0.4271
\end{verbatim}
\end{kframe}
\end{knitrout}

\subsubsection{Bulk-only special case (no GPD tail)}
\label{subsec:overview-nogpd}

Setting \code{GPD = FALSE} removes the splicing step so that
$f(y\mid x;\Theta)\equiv f_{\mathrm{DP}}(y\mid x;\Theta)$, and
$L(\Theta,\Phi;\mathcal{D})$ reduces to the bulk-mixture likelihood in
$\Theta$. All estimation and prediction targets described below remain
available, with tail-specific terms omitted. The wrapper \code{dpmix()}
fits this bulk-only model, while \code{dpmgpd()} fits the augmented
(spliced) model.

\subsection{Posterior functionals, interval summaries and prediction}
\label{subsec:overview-estimation}

The estimands reported by \pkg{CausalMixGPD} are functionals of the
estimated conditional distribution $F(\cdot\mid x)$ from
Section~\ref{sec:prelim}. In this article we focus on four conditional
functionals: density $f(y\mid x)$, survival $S(y\mid x)=1-F(y\mid x)$,
quantiles $Q(\tau\mid x)$, and mean-type summaries.

Let $(\Theta^{(s)},\Phi^{(s)})_{s=1}^S$ denote post-burn-in draws from
the posterior. For density, survival, and quantiles, \pkg{CausalMixGPD}
evaluates the corresponding draw-wise functional using the spliced
formulas \eqref{eq:splice_cdf}--\eqref{eq:splice_quantile_case} (with
the truncated bulk mixture), and then summarizes across draws.
Specifically, for a generic scalar functional $T(\Theta,\Phi;x)$ (e.g.,
$f(y_0\mid x)$ or $Q(\tau\mid x)$), $$
  \widehat{T}(x)=\frac{1}{S}\sum_{s=1}^S T(\Theta^{(s)},\Phi^{(s)};x),
$$ and uncertainty is summarized either by equal-tailed credible
intervals (posterior quantiles) or by highest posterior density (HPD)
intervals (shortest intervals covering the target posterior mass).
Equal-tailed intervals and general Bayesian posterior summaries follow
standard practice \citep{Gelman2013BDA}; HPD intervals are computed from
the MCMC draws using established implementations \citep{coda_pkg}.

Mean-type summaries are handled differently in the package. For
\code{type = "mean"} and \code{type = "rmean"}, \pkg{CausalMixGPD} uses
posterior predictive Monte Carlo within each posterior draw rather than
direct plug-in evaluation. For a fixed covariate value $x$, within draw
$s$ the package generates $$
  y_{1}^{(s)},\ldots,y_{M}^{(s)} \;\sim\; f(\cdot\mid x;\Theta^{(s)},\Phi^{(s)}),
$$ and computes $$
  \mu^{(s)}(x)=\frac{1}{M}\sum_{m=1}^M y_{m}^{(s)},
  \qquad
  \mu_{c}^{(s)}(x)=\frac{1}{M}\sum_{m=1}^M \min\{y_{m}^{(s)},c\},
$$ corresponding to the conditional mean and the restricted mean at
cutoff $c$, respectively. Posterior summaries and intervals are then
formed from $\{\mu^{(s)}(x)\}_{s=1}^S$ or $\{\mu_c^{(s)}(x)\}_{s=1}^S$.
Under the spliced model, the conditional mean is infinite when
$\xi\ge 1$; accordingly, \pkg{CausalMixGPD} flags \code{type = "mean"}
as infinite when there is posterior mass with $\xi\ge 1$, while
\code{type = "rmean"} remains finite by construction.

Using the same fitted model \code{fit} from
Section~\ref{subsec:overview-model}, we now compute posterior summaries
on the original training covariates (default behavior of
\code{predict()} when \code{newdata} is not supplied).

\begin{knitrout}\footnotesize
\definecolor{shadecolor}{rgb}{0.969, 0.969, 0.969}\color{fgcolor}\begin{kframe}
\begin{alltt}
\hldef{q_grid} \hlkwb{<-} \hlkwd{c}\hldef{(}\hlnum{0.50}\hldef{,} \hlnum{0.90}\hldef{,} \hlnum{0.95}\hldef{,} \hlnum{0.99}\hldef{)}
\hldef{x_eval} \hlkwb{<-} \hldef{dat[}\hlnum{1}\hlopt{:}\hlnum{20}\hldef{,} \hlkwd{c}\hldef{(}\hlsng{"x1"}\hldef{,} \hlsng{"x2"}\hldef{,} \hlsng{"x3"}\hldef{)]}

\hldef{est_quant} \hlkwb{<-} \hlkwd{predict}\hldef{(fit,} \hlkwc{newdata} \hldef{= x_eval,} \hlkwc{type} \hldef{=} \hlsng{"quantile"}\hldef{,} \hlkwc{index} \hldef{= q_grid,}
                     \hlkwc{interval} \hldef{=} \hlsng{"hpd"}\hldef{,} \hlkwc{level} \hldef{=} \hlnum{0.95}\hldef{)}
\end{alltt}
\end{kframe}
\end{knitrout}



\begin{knitrout}\footnotesize
\definecolor{shadecolor}{rgb}{0.969, 0.969, 0.969}\color{fgcolor}\begin{kframe}
\begin{verbatim}
   id index   estimate      lower     upper
1   1  0.50  1.8210682 0.50229729  2.909784
2   1  0.90  4.6807704 3.19430382  6.411481
3   1  0.95  6.2968609 4.03365615  9.313934
4   1  0.99 12.7849025 7.27892012 23.359469
5   2  0.50  0.7598295 0.17676829  1.952736
6   2  0.90  3.5837439 2.50637730  5.427703
7   2  0.95  4.7603592 3.28801838  7.299773
8   2  0.99  9.3660456 5.52198497 14.274001
9   3  0.50  1.0595683 0.03923384  2.729271
10  3  0.90  3.5216020 1.16265972  5.846167
\end{verbatim}
\end{kframe}
\end{knitrout}

For a new covariate value $x^\star$, prediction is based on the
posterior predictive distribution $$
  p(y^\star\mid x^\star,\mathcal{D})
  =\int f(y^\star\mid x^\star;\Theta,\Phi)\,p(\Theta,\Phi\mid \mathcal{D})\,d\Theta\,d\Phi.
$$ The \code{predict()} method returns posterior summaries for the same
functional targets as in Section~\ref{subsec:overview-estimation}. For
\code{type = "density"} and \code{type = "survival"}, the user supplies
a grid of $y$ values and the method returns pointwise posterior means
with interval bounds computed draw-by-draw. For
\code{type = "quantile"}, the user supplies quantile indices $\tau$ and
the method returns posterior summaries of $Q(\tau\mid x^\star)$ computed
using the spliced quantile representation in
\eqref{eq:splice_quantile_case}. For \code{type = "mean"} and
\code{type = "rmean"}, prediction uses posterior predictive Monte Carlo
as described above, with $x$ replaced by $x^\star$ and with
user-controlled $M$ through \code{nsim\_mean}.

\begin{knitrout}\footnotesize
\definecolor{shadecolor}{rgb}{0.969, 0.969, 0.969}\color{fgcolor}\begin{kframe}
\begin{alltt}
\hldef{q_levels} \hlkwb{<-} \hlkwd{c}\hldef{(}\hlnum{0.25}\hldef{,} \hlnum{0.50}\hldef{,} \hlnum{0.75}\hldef{)}
\hldef{x_new} \hlkwb{<-} \hlkwd{as.data.frame}\hldef{(}\hlkwd{lapply}\hldef{(}
  \hldef{dat[,} \hlkwd{c}\hldef{(}\hlsng{"x1"}\hldef{,} \hlsng{"x2"}\hldef{,} \hlsng{"x3"}\hldef{)],}
  \hldef{quantile,}
  \hlkwc{probs} \hldef{= q_levels,}
  \hlkwc{na.rm} \hldef{=} \hlnum{TRUE}
\hldef{))}
\hlkwd{rownames}\hldef{(x_new)} \hlkwb{<-} \hlkwd{c}\hldef{(}\hlsng{"q25"}\hldef{,} \hlsng{"q50"}\hldef{,} \hlsng{"q75"}\hldef{)}
\hldef{y_grid} \hlkwb{<-} \hlkwd{seq}\hldef{(}\hlnum{0}\hldef{,} \hlnum{10}\hldef{,} \hlkwc{length.out} \hldef{=} \hlnum{200}\hldef{)}
\hldef{p_grid} \hlkwb{<-} \hlkwd{c}\hldef{(}\hlnum{0.50}\hldef{,} \hlnum{0.90}\hldef{,} \hlnum{0.95}\hldef{,} \hlnum{0.99}\hldef{)}
\end{alltt}
\end{kframe}
\end{knitrout}

\begin{knitrout}\footnotesize
\definecolor{shadecolor}{rgb}{0.969, 0.969, 0.969}\color{fgcolor}\begin{kframe}
\begin{alltt}
\hldef{pdens} \hlkwb{<-} \hlkwd{predict}\hldef{(fit,} \hlkwc{newdata} \hldef{= x_new,} \hlkwc{y} \hldef{= y_grid,} \hlkwc{type} \hldef{=} \hlsng{"density"}\hldef{,}
                 \hlkwc{interval} \hldef{=} \hlsng{"credible"}\hldef{,} \hlkwc{level} \hldef{=} \hlnum{0.95}\hldef{)}
\end{alltt}
\end{kframe}
\end{knitrout}

\begin{knitrout}\footnotesize
\definecolor{shadecolor}{rgb}{0.969, 0.969, 0.969}\color{fgcolor}\begin{kframe}
\begin{alltt}
\hldef{psurv} \hlkwb{<-} \hlkwd{predict}\hldef{(fit,} \hlkwc{newdata} \hldef{= x_new,} \hlkwc{y} \hldef{= y_grid,} \hlkwc{type} \hldef{=} \hlsng{"survival"}\hldef{,}
                 \hlkwc{interval} \hldef{=} \hlsng{"credible"}\hldef{,} \hlkwc{level} \hldef{=} \hlnum{0.95}\hldef{)}
\end{alltt}
\end{kframe}
\end{knitrout}

\begin{knitrout}\footnotesize
\definecolor{shadecolor}{rgb}{0.969, 0.969, 0.969}\color{fgcolor}\begin{kframe}
\begin{alltt}
\hldef{pquant} \hlkwb{<-} \hlkwd{predict}\hldef{(fit,} \hlkwc{newdata} \hldef{= x_new,} \hlkwc{type} \hldef{=} \hlsng{"quantile"}\hldef{,} \hlkwc{index} \hldef{= p_grid,}
                  \hlkwc{interval} \hldef{=} \hlsng{"hpd"}\hldef{,} \hlkwc{level} \hldef{=} \hlnum{0.95}\hldef{)}
\end{alltt}
\end{kframe}
\end{knitrout}



\begin{knitrout}\footnotesize
\definecolor{shadecolor}{rgb}{0.969, 0.969, 0.969}\color{fgcolor}\begin{kframe}
\begin{verbatim}
             x1          x2            x3
q25 -0.59851252 -0.57151934 -0.7224762264
q50  0.02874559 -0.01658997  0.0004347721
q75  0.69727001  0.41546653  0.5516024827
 [1] 0.00000000 0.05025126 0.10050251 0.15075377 0.20100503
 [6] 0.25125628 0.30150754 0.35175879 0.40201005 0.45226131
[1] 0.50 0.90 0.95 0.99
\end{verbatim}
\end{kframe}
\end{knitrout}

\begin{knitrout}\footnotesize
\definecolor{shadecolor}{rgb}{0.969, 0.969, 0.969}\color{fgcolor}\begin{kframe}
\begin{verbatim}
   id          y       density         lower         upper
1   1 0.00000000 1.000000e-300 1.000000e-300 1.000000e-300
2   1 0.05025126  1.050943e+00  4.391043e-01  2.415552e+00
3   1 0.10050251  7.648774e-01  3.187806e-01  1.422371e+00
4   1 0.15075377  6.226713e-01  2.719930e-01  1.181847e+00
5   1 0.20100503  5.314829e-01  2.231017e-01  9.747405e-01
6   1 0.25125628  4.720142e-01  1.843810e-01  8.605983e-01
7   1 0.30150754  4.276837e-01  1.595827e-01  7.921671e-01
8   1 0.35175879  3.923244e-01  1.433871e-01  7.182820e-01
9   1 0.40201005  3.629111e-01  1.324248e-01  6.402643e-01
10  1 0.45226131  3.378112e-01  1.223972e-01  5.899420e-01
\end{verbatim}
\end{kframe}
\end{knitrout}

\begin{knitrout}\footnotesize
\definecolor{shadecolor}{rgb}{0.969, 0.969, 0.969}\color{fgcolor}\begin{kframe}
\begin{verbatim}
   id          y  survival     lower     upper
1   1 0.00000000 1.0000000 1.0000000 1.0000000
2   1 0.05025126 0.9055000 0.8106256 0.9673990
3   1 0.10050251 0.8613613 0.7228659 0.9408555
4   1 0.15075377 0.8266686 0.6671500 0.9196784
5   1 0.20100503 0.7978666 0.6270789 0.9013449
6   1 0.25125628 0.7727374 0.5958568 0.8848708
7   1 0.30150754 0.7501789 0.5701748 0.8714609
8   1 0.35175879 0.7296054 0.5482139 0.8605438
9   1 0.40201005 0.7106504 0.5288923 0.8490909
10  1 0.45226131 0.6930600 0.5084070 0.8356635
\end{verbatim}
\end{kframe}
\end{knitrout}

\begin{knitrout}\footnotesize
\definecolor{shadecolor}{rgb}{0.969, 0.969, 0.969}\color{fgcolor}\begin{kframe}
\begin{verbatim}
   id index  estimate     lower     upper
1   1  0.50 1.4823780 0.4926661  2.711455
2   1  0.90 3.9245522 2.9765853  5.409868
3   1  0.95 5.0905154 3.6744175  7.986932
4   1  0.99 9.7925408 6.0698483 20.129287
5   2  0.50 1.0055660 0.9595097  1.065311
6   2  0.90 3.6805966 3.3302733  4.009358
7   2  0.95 4.8939867 4.3663054  5.726796
8   2  0.99 9.7018133 7.5487018 13.250233
9   3  0.50 0.9570955 0.1153466  2.018721
10  3  0.90 4.3029306 3.1282140  5.638944
\end{verbatim}
\end{kframe}
\end{knitrout}

\subsection{Causal extension and propensity score augmentation}
\label{subsec:overview-causal}

This package extends the one-arm conditional mixture regression
framework to two-arm causal studies by fitting an outcome model under
treatment and under control, optionally augmenting the outcome
covariates with a PS summary. Throughout this section
we write the model in its default \emph{augmented spliced} form (bulk
mixture + GPD tail), using the same conditional distributional building
blocks introduced earlier (so we do not restate the spliced density; see
\eqref{eq:splice_pdf} and the corresponding conditional functionals
defined in the previous sections).

\subsubsection{Data, model components, and likelihood factorization}

Let $$
  \mathcal{D}=\{(y_i,a_i,\boldsymbol{x}_i)\}_{i=1}^n,
  \qquad a_i\in\{0,1\},
$$ where $a_i=1$ denotes treatment and $a_i=0$ denotes control. The
causal interface consists of two pieces.

First, when PS augmentation is enabled, a binary regression model is
used for the treatment mechanism \begin{align}
  a_i \mid \boldsymbol{x}_i, \boldsymbol{\gamma}
  &\sim \mathrm{Bernoulli}\big(e_{\boldsymbol{\gamma}}(\boldsymbol{x}_i)\big),
  \label{eq:overview_ps_part1}\\
  h\!\left(e_{\boldsymbol{\gamma}}(\boldsymbol{x}_i)\right)
  &= \gamma_0 + \boldsymbol{x}_i^\top\boldsymbol{\gamma},
  \qquad h\in\{\text{logit},\text{probit}\},
  \label{eq:overview_ps_part2}
\end{align} which is consistent with the definition and balancing
property in \eqref{eq:ps_def}--\eqref{eq:ps_balance}. The package also
supports a \texttt{"naive"} option, in which $e(\boldsymbol{x})$ is
taken to be constant (no covariate dependence).

Second, the outcome model is fit separately within each arm. Denote by
$F_a(\,\cdot\mid\boldsymbol{r})$ the fitted conditional outcome
distribution under arm $a\in\{0,1\}$, where $\boldsymbol{r}$ is the
covariate vector used by the outcome model. With PS augmentation, the
package constructs an \emph{augmented} covariate vector \begin{equation}
  \boldsymbol{r}(\boldsymbol{x})
  =
  \Big(1,\ \boldsymbol{x}^\top,\ \psi\big(\widehat e(\boldsymbol{x})\big)\Big)^\top,
  \label{eq:overview_augmented_covariate}
\end{equation} where $\widehat e(\boldsymbol{x})$ is a posterior summary
of the PS (posterior mean or median, controlled by
\texttt{ps\_summary}), and $\psi(\cdot)$ is either
$\psi(p)=\text{logit}(p)$ or $\psi(p)=p$ (controlled by
\texttt{ps\_scale}). When PS augmentation is disabled,
$\boldsymbol{r}(\boldsymbol{x})=(1,\boldsymbol{x}^\top)^\top$.

Within each arm, the conditional density $f_a(y\mid \boldsymbol{r})$
follows the same mixture regression + optional GPD tail specification
described in the one-arm section (kernel choice, link structure
$g(\cdot)$ for covariate-dependent parameters, and either \texttt{"sb"}
or \texttt{"crp"} backend for the DPM). Let
$\boldsymbol{\Theta}_a$ collect the arm-specific mixture and tail
parameters. The arm-wise likelihood contributions are $$
  L_a(\boldsymbol{\Theta}_a;\mathcal{D})
  =
  \prod_{i:\,a_i=a} f_a\!\left(y_i \mid \boldsymbol{r}(\boldsymbol{x}_i); \boldsymbol{\Theta}_a\right),
  \qquad a\in\{0,1\},
$$ with $f_a(\cdot\mid\cdot)$ given by the spliced construction already
defined (bulk mixture below threshold and GPD exceedances above
threshold). Practically, the implementation uses a two-stage workflow:
(i) fit the PS model using
\eqref{eq:overview_ps_part1}--\eqref{eq:overview_ps_part2}, then (ii)
compute $\widehat e(\boldsymbol{x}_i)$ and build
$\boldsymbol{r}(\boldsymbol{x}_i)$ in
\eqref{eq:overview_augmented_covariate} before fitting the two
arm-specific outcome models.

\subsubsection{Prior specification (PS then outcome)}

For the PS regression coefficients, the default prior is independent
Normal, $$
  \gamma_0 \sim \mathcal{N}(m_0,s_0^2),
  \qquad
  \gamma_k \sim \mathcal{N}(m_0,s_0^2),
$$ with defaults controlled by \texttt{ps\_prior}.

For each arm's outcome model, priors follow the same structure as the
one-arm case: a DPM prior (with either truncated stick-breaking
or CRP allocation, implemented in \texttt{nimble}
\citep{devalpine2017nimble,nimble_pkg}) and kernel-/tail-specific priors
on the bulk and GPD parameters. Since the causal fit is simply a pair of
one-arm fits (treated and control), all customization options for
kernels, links, and tail inclusion apply arm-wise.

\subsubsection{MCMC algorithm, diagnostics, and parameter extraction}

The causal fit runs MCMC in two stages, each with its own iteration
controls.

(i) PS stage: posterior simulation for $\boldsymbol{\gamma}$ in
    \eqref{eq:overview_ps_part1}--\eqref{eq:overview_ps_part2} using
    \texttt{mcmc\_ps} settings. The PS posterior draws are summarized to
    obtain $\widehat e(\boldsymbol{x}_i)$, which is optionally
    transformed via \texttt{ps\_scale} and appended to the outcome
    covariates as in \eqref{eq:overview_augmented_covariate}.

(ii) Outcome stage: arm-specific posterior simulation for
     $\boldsymbol{\Theta}_0$ and $\boldsymbol{\Theta}_1$ using
     \texttt{mcmc\_outcome} settings, with either \texttt{"sb"} or
     \texttt{"crp"} backend.

For routine convergence checks, users can inspect trace plots and
autocorrelation for monitored parameters and (when running multiple
chains) compare between-chain and within-chain mixing. The package
stores MCMC samples in the fitted object, and standard diagnostics can
be applied after converting samples to \texttt{coda} objects (e.g.,
effective sample size, Geweke diagnostics, and HPD intervals via
\texttt{coda::HPDinterval}; see the \texttt{coda} reference at the end
of this message). Posterior parameter summaries can be extracted via the
exported \texttt{params()} method, which returns arm-specific parameter
tables for causal fits.

\subsubsection{Estimation and interval summaries for causal contrasts}

All reported causal quantities are computed as functionals of the fitted
arm-specific conditional distributions. Let $T_a(\boldsymbol{x})$ denote
any conditional functional under arm $a$ (density, survival, quantile,
or mean/restricted mean), computed from
$F_a(\cdot\mid\boldsymbol{r}(\boldsymbol{x}))$ using the same
definitions as in the one-arm case. The corresponding conditional causal
contrast is $$
  \Delta_T(\boldsymbol{x}) = T_1(\boldsymbol{x}) - T_0(\boldsymbol{x}).
$$ The package computes $\Delta_T(\boldsymbol{x})$ \emph{draw-by-draw}
from posterior samples, then summarizes $\Delta_T(\boldsymbol{x})$ by a
posterior mean/median and uncertainty intervals. Equal-tailed credible
intervals are formed by posterior quantiles, while HPD intervals are
computed using \texttt{coda::HPDinterval}.

The main causal estimators provided by the package are:

(a) Conditional effects at covariate profiles. For a user-supplied set
    of covariate rows $\{\boldsymbol{x}^\star_m\}_{m=1}^M$, $$
      \mathrm{CATE}(\boldsymbol{x}^\star_m)=\mu_1(\boldsymbol{x}^\star_m)-\mu_0(\boldsymbol{x}^\star_m),
      \qquad
      \mathrm{CQTE}(\tau\mid \boldsymbol{x}^\star_m)
      =
      Q_1(\tau\mid \boldsymbol{x}^\star_m)-Q_0(\tau\mid \boldsymbol{x}^\star_m),
    $$ where
    $\mu_a(\boldsymbol{x})=\mathbb{E}(Y\mid A=a,\boldsymbol{X}=\boldsymbol{x})$
    and $Q_a(\tau\mid\boldsymbol{x})$ is the conditional quantile
    function. In heavy-tailed settings, $\mu_a(\boldsymbol{x})$ can be
    unstable when the implied tail shape yields infinite mean; the
    package therefore supports both \texttt{type="mean"} and
    \texttt{type="rmean"} for (restricted) mean-based effects, with
    restricted means computed as
    $\mathbb{E}[\min\{Y,c\}\mid A=a,\boldsymbol{X}=\boldsymbol{x}]$ at
    user-specified cutoff $c$.

(b) Marginal (standardized) summaries. The package reports marginal ATE
    and marginal QTE-style summaries by empirical standardization over
    covariate profiles: \begin{align}
      \mathrm{ATE}
      &= \frac{1}{n}\sum_{i=1}^n \Big\{\mu_1(\boldsymbol{x}_i)-\mu_0(\boldsymbol{x}_i)\Big\},
      \label{eq:overview_ate_standardized}\\
      \mathrm{QTE}(\tau)
      &= Q_1^{\mathrm{m}}(\tau)-Q_0^{\mathrm{m}}(\tau),\qquad
      Q_a^{\mathrm{m}}(\tau)=\inf\{y:\,F_a^{\mathrm{m}}(y)\ge \tau\},\quad
      F_a^{\mathrm{m}}(y)=\frac{1}{n}\sum_{i=1}^n F_a(y\mid \boldsymbol{x}_i).
      \label{eq:overview_qte_standardized}
    \end{align} Treated-standardized analogues (ATT, QTT) are obtained
    by replacing the full-sample covariate empirical distribution with
    the treated-subsample empirical distribution $\{i:a_i=1\}$. In the
    current implementation, \texttt{ate()} and \texttt{att()} aggregate
    draw-wise conditional mean contrasts across rows, while
    \texttt{qte()} and \texttt{qtt()} compute draw-wise arm-specific
    marginal quantiles and then contrast them.

\subsubsection{Posterior prediction at new covariate profiles}

Prediction in the causal setup proceeds by producing arm-specific
predictive summaries at new covariates and then differencing when
required. For new covariates $\boldsymbol{x}^\star$, the package first
predicts the PS (if enabled), forms
$\boldsymbol{r}(\boldsymbol{x}^\star)$ via
\eqref{eq:overview_augmented_covariate}, and then evaluates the desired
distributional functional under both arms. For \texttt{type = "mean"}
and \texttt{type = "quantile"}, \texttt{predict()} returns
treated-minus-control summaries and stores arm-specific predictions as
attributes; for \texttt{type = "density"}, \texttt{"survival"}, and
\texttt{"prob"}, it returns arm-specific summaries directly.

A recommended way to build an interpretable prediction grid is to use
coordinate-wise empirical quartiles. Let $q_{k,p}$ be the $p$th sample
quantile of covariate $x_k$ for $p\in\{0.25,0.5,0.75\}$. The grid $$
  \mathcal{X}_\star = \prod_{k=1}^p \{q_{k,0.25},q_{k,0.5},q_{k,0.75}\}
$$ yields $3^p$ representative covariate profiles at which
$\mathrm{CATE}$, $\mathrm{CQTE}$, survival contrasts, and other
predictive contrasts can be reported.

\subsubsection{User-facing workflow (code)}

The following example uses the formula interface and the high-level
wrapper, fitting the augmented spliced model by default.

\begin{knitrout}\footnotesize
\definecolor{shadecolor}{rgb}{0.969, 0.969, 0.969}\color{fgcolor}\begin{kframe}
\begin{alltt}
\hlkwd{data}\hldef{(}\hlsng{"causal_pos500_p3_k2"}\hldef{,} \hlkwc{package} \hldef{=} \hlsng{"CausalMixGPD"}\hldef{)}
\hldef{dat} \hlkwb{<-} \hldef{causal_pos500_p3_k2}
\hldef{df}  \hlkwb{<-} \hlkwd{data.frame}\hldef{(}\hlkwc{y} \hldef{= dat}\hlopt{$}\hldef{y,} \hlkwc{A} \hldef{= dat}\hlopt{$}\hldef{A, dat}\hlopt{$}\hldef{X)}
\end{alltt}
\end{kframe}
\end{knitrout}

\begin{knitrout}\footnotesize
\definecolor{shadecolor}{rgb}{0.969, 0.969, 0.969}\color{fgcolor}\begin{kframe}
\begin{alltt}
\hldef{mcmc_fixed} \hlkwb{<-} \hlkwd{list}\hldef{(}
  \hlkwc{niter} \hldef{=} \hlnum{200}\hldef{,}
  \hlkwc{nburnin} \hldef{=} \hlnum{50}\hldef{,}
  \hlkwc{thin} \hldef{=} \hlnum{1}\hldef{,}
  \hlkwc{nchains} \hldef{=} \hlnum{1}\hldef{,}
  \hlkwc{seed} \hldef{=} \hlnum{1}
\hldef{)}
\end{alltt}
\end{kframe}
\end{knitrout}

\begin{knitrout}\footnotesize
\definecolor{shadecolor}{rgb}{0.969, 0.969, 0.969}\color{fgcolor}\begin{kframe}
\begin{alltt}
\hldef{cfit} \hlkwb{<-} \hlkwd{dpmgpd}\hldef{(}
  \hlkwc{formula} \hldef{= y} \hlopt{~} \hldef{x1} \hlopt{+} \hldef{x2} \hlopt{+} \hldef{x3,}
  \hlkwc{data} \hldef{= df,}
  \hlkwc{treat} \hldef{=} \hlsng{"A"}\hldef{,}
  \hlkwc{backend} \hldef{=} \hlsng{"sb"}\hldef{,}
  \hlkwc{kernel} \hldef{=} \hlsng{"lognormal"}\hldef{,}
  \hlkwc{components} \hldef{=} \hlnum{5}\hldef{,}
  \hlkwc{PS} \hldef{=} \hlsng{"logit"}\hldef{,}
  \hlkwc{ps_scale} \hldef{=} \hlsng{"logit"}\hldef{,}
  \hlkwc{ps_summary} \hldef{=} \hlsng{"mean"}\hldef{,}
  \hlkwc{mcmc_outcome} \hldef{= mcmc_fixed,}
  \hlkwc{mcmc_ps} \hldef{= mcmc_fixed}
\hldef{)}
\end{alltt}
\end{kframe}
\end{knitrout}



\begin{knitrout}\footnotesize
\definecolor{shadecolor}{rgb}{0.969, 0.969, 0.969}\color{fgcolor}\begin{kframe}
\begin{verbatim}
CausalMixGPD causal fit
PS model: Bayesian logit (A | X) 
Outcome (treated): backend = sb | kernel = lognormal 
Outcome (control): backend = sb | kernel = lognormal 
GPD tail (treated/control): TRUE / TRUE 
\end{verbatim}
\end{kframe}
\end{knitrout}

\begin{knitrout}\footnotesize
\definecolor{shadecolor}{rgb}{0.969, 0.969, 0.969}\color{fgcolor}\begin{kframe}
\begin{alltt}
\hldef{ate_hat} \hlkwb{<-} \hlkwd{ate}\hldef{(cfit,} \hlkwc{level} \hldef{=} \hlnum{0.90}\hldef{,} \hlkwc{interval} \hldef{=} \hlsng{"credible"}\hldef{)}
\end{alltt}
\end{kframe}
\end{knitrout}

\begin{knitrout}\footnotesize
\definecolor{shadecolor}{rgb}{0.969, 0.969, 0.969}\color{fgcolor}\begin{kframe}
\begin{verbatim}
  id estimate     lower    upper
1  1 133.6313 0.5208869 789.4111
\end{verbatim}
\end{kframe}
\end{knitrout}

\begin{knitrout}\footnotesize
\definecolor{shadecolor}{rgb}{0.969, 0.969, 0.969}\color{fgcolor}\begin{kframe}
\begin{alltt}
\hldef{qte_hat} \hlkwb{<-} \hlkwd{qte}\hldef{(cfit,} \hlkwc{probs} \hldef{=} \hlkwd{c}\hldef{(}\hlnum{0.50}\hldef{,} \hlnum{0.90}\hldef{,} \hlnum{0.95}\hldef{),} \hlkwc{level} \hldef{=} \hlnum{0.90}\hldef{,} \hlkwc{interval} \hldef{=} \hlsng{"hpd"}\hldef{)}
\end{alltt}
\end{kframe}
\end{knitrout}

\begin{knitrout}\footnotesize
\definecolor{shadecolor}{rgb}{0.969, 0.969, 0.969}\color{fgcolor}\begin{kframe}
\begin{verbatim}
  id index  estimate       lower     upper
1  1  0.50 -0.230319 -0.72616380 0.2434466
2  1  0.90  1.581025  0.04824207 2.8446616
3  1  0.95  3.424051 -1.11525657 7.0265043
\end{verbatim}
\end{kframe}
\end{knitrout}

\begin{knitrout}\footnotesize
\definecolor{shadecolor}{rgb}{0.969, 0.969, 0.969}\color{fgcolor}\begin{kframe}
\begin{alltt}
\hldef{Xnew} \hlkwb{<-} \hldef{df[}\hlnum{1}\hlopt{:}\hlnum{5}\hldef{,} \hlkwd{c}\hldef{(}\hlsng{"x1"}\hldef{,}\hlsng{"x2"}\hldef{,}\hlsng{"x3"}\hldef{)]}
\end{alltt}
\end{kframe}
\end{knitrout}

\begin{knitrout}\footnotesize
\definecolor{shadecolor}{rgb}{0.969, 0.969, 0.969}\color{fgcolor}\begin{kframe}
\begin{alltt}
\hldef{cate_hat} \hlkwb{<-} \hlkwd{cate}\hldef{(cfit,} \hlkwc{newdata} \hldef{= Xnew,} \hlkwc{type} \hldef{=} \hlsng{"mean"}\hldef{,} \hlkwc{level} \hldef{=} \hlnum{0.90}\hldef{)}
\end{alltt}
\end{kframe}
\end{knitrout}

\begin{knitrout}\footnotesize
\definecolor{shadecolor}{rgb}{0.969, 0.969, 0.969}\color{fgcolor}\begin{kframe}
\begin{verbatim}
  id    estimate      lower     upper
1  1  0.15628140 -0.8192142 1.2088089
2  2 -0.01701753 -1.8635060 2.9794351
3  3  0.80363038 -0.2959428 2.7214028
4  4 -0.19121860 -1.7339494 2.1439422
5  5  0.10625831 -0.1619069 0.6342523
\end{verbatim}
\end{kframe}
\end{knitrout}

\begin{knitrout}\footnotesize
\definecolor{shadecolor}{rgb}{0.969, 0.969, 0.969}\color{fgcolor}\begin{kframe}
\begin{alltt}
\hldef{cqte_hat} \hlkwb{<-} \hlkwd{cqte}\hldef{(cfit,} \hlkwc{newdata} \hldef{= Xnew,} \hlkwc{probs} \hldef{=} \hlkwd{c}\hldef{(}\hlnum{0.50}\hldef{,} \hlnum{0.90}\hldef{),} \hlkwc{level} \hldef{=} \hlnum{0.90}\hldef{)}
\end{alltt}
\end{kframe}
\end{knitrout}

\begin{knitrout}\footnotesize
\definecolor{shadecolor}{rgb}{0.969, 0.969, 0.969}\color{fgcolor}\begin{kframe}
\begin{verbatim}
   id index    estimate       lower     upper
1   1   0.5  0.09848583 -0.32638030 0.5465127
2   1   0.9  0.14929438 -0.92614571 1.4969607
3   2   0.5 -0.03731305 -0.42919076 0.3870325
4   2   0.9 -0.14538125 -1.84931903 2.0484655
5   3   0.5  0.37020524 -0.12156714 0.8630646
6   3   0.9  1.03612061 -0.43364329 3.0147317
7   4   0.5 -0.06317799 -0.47837124 0.2813245
8   4   0.9 -0.23852380 -2.04943975 2.6413435
9   5   0.5  0.04424529 -0.08159786 0.2133258
10  5   0.9  0.22121881 -0.23041000 1.2275466
\end{verbatim}
\end{kframe}
\end{knitrout}

\begin{knitrout}\footnotesize
\definecolor{shadecolor}{rgb}{0.969, 0.969, 0.969}\color{fgcolor}\begin{kframe}
\begin{alltt}
\hldef{qs}    \hlkwb{<-} \hlkwd{c}\hldef{(}\hlnum{0.25}\hldef{,} \hlnum{0.50}\hldef{,} \hlnum{0.75}\hldef{)}
\hldef{Xgrid} \hlkwb{<-} \hlkwd{expand.grid}\hldef{(}\hlkwd{lapply}\hldef{(df[}\hlkwd{c}\hldef{(}\hlsng{"x1"}\hldef{,}\hlsng{"x2"}\hldef{,}\hlsng{"x3"}\hldef{)], quantile,} \hlkwc{probs} \hldef{= qs))}
\end{alltt}
\end{kframe}
\end{knitrout}

\begin{knitrout}\footnotesize
\definecolor{shadecolor}{rgb}{0.969, 0.969, 0.969}\color{fgcolor}\begin{kframe}
\begin{alltt}
\hldef{pred_q} \hlkwb{<-} \hlkwd{predict}\hldef{(cfit,} \hlkwc{newdata} \hldef{= Xgrid,} \hlkwc{type} \hldef{=} \hlsng{"quantile"}\hldef{,} \hlkwc{p} \hldef{=} \hlkwd{c}\hldef{(}\hlnum{0.50}\hldef{,} \hlnum{0.90}\hldef{))}
\end{alltt}
\end{kframe}
\end{knitrout}

\begin{knitrout}\footnotesize
\definecolor{shadecolor}{rgb}{0.969, 0.969, 0.969}\color{fgcolor}\begin{kframe}
\begin{verbatim}
data frame with 0 columns and 0 rows
\end{verbatim}
\end{kframe}
\end{knitrout}

\begin{knitrout}\footnotesize
\definecolor{shadecolor}{rgb}{0.969, 0.969, 0.969}\color{fgcolor}\begin{kframe}
\begin{alltt}
\hldef{pred_s} \hlkwb{<-} \hlkwd{predict}\hldef{(cfit,} \hlkwc{newdata} \hldef{= Xgrid,} \hlkwc{type} \hldef{=} \hlsng{"survival"}\hldef{,} \hlkwc{y} \hldef{=} \hlkwd{rep}\hldef{(}\hlnum{4}\hldef{,} \hlkwd{nrow}\hldef{(Xgrid)))}
\end{alltt}
\end{kframe}
\end{knitrout}

\begin{knitrout}\footnotesize
\definecolor{shadecolor}{rgb}{0.969, 0.969, 0.969}\color{fgcolor}\begin{kframe}
\begin{verbatim}
data frame with 0 columns and 0 rows
\end{verbatim}
\end{kframe}
\end{knitrout}

\begin{knitrout}\footnotesize
\definecolor{shadecolor}{rgb}{0.969, 0.969, 0.969}\color{fgcolor}\begin{kframe}
\begin{alltt}
\hldef{pred_q_trt} \hlkwb{<-} \hlkwd{attr}\hldef{(pred_q,} \hlsng{"trt"}\hldef{)}
\hldef{pred_q_con} \hlkwb{<-} \hlkwd{attr}\hldef{(pred_q,} \hlsng{"con"}\hldef{)}
\end{alltt}
\end{kframe}
\end{knitrout}

\begin{knitrout}\footnotesize
\definecolor{shadecolor}{rgb}{0.969, 0.969, 0.969}\color{fgcolor}\begin{kframe}
\begin{verbatim}
   id index  estimate      lower     upper
1   1   0.5 1.0871813 0.54652435 2.3651155
2   1   0.9 3.1515644 1.79120976 4.8641579
3   2   0.5 0.4514065 0.20455294 0.8192195
4   2   0.9 1.6172941 0.90073772 2.9213755
5   3   0.5 0.2518916 0.06924277 0.5237770
6   3   0.9 1.1367718 0.52731902 3.1330499
7   4   0.5 1.9931938 1.27091630 3.0588945
8   4   0.9 4.4549007 3.86971209 5.8554098
9   5   0.5 0.7953029 0.44198724 1.3200333
10  5   0.9 2.5894862 1.80677822 3.6550876
   id index  estimate      lower     upper
1   1   0.5 1.5998935 0.97997952 2.0270002
2   1   0.9 3.5348458 2.92201489 4.3121050
3   2   0.5 0.4242338 0.26661400 0.7774886
4   2   0.9 1.4014587 0.98846982 2.2199421
5   3   0.5 0.1427877 0.06161788 0.3169749
6   3   0.9 0.5103574 0.24296364 0.9483756
7   4   0.5 2.0936708 1.84122405 2.3247500
8   4   0.9 4.1776855 3.76068895 4.8194525
9   5   0.5 0.7199800 0.50469483 1.0582684
10  5   0.9 2.0859245 1.62457363 2.7495954
\end{verbatim}
\end{kframe}
\end{knitrout}

\subsubsection{Bulk-only causal mixtures (no GPD tail)}

To run the same two-arm regression without a GPD tail, use the bulk-only
wrapper \texttt{dpmix()} (or set \texttt{GPD=FALSE} at bundle
construction). All causal estimand and prediction functions
(\texttt{ate}, \texttt{qte}, \texttt{cate}, \texttt{cqte}, and
\texttt{predict}) operate the same way, now interpreting
$F_a(\cdot\mid\boldsymbol{x})$ as the bulk mixture distribution.

\section{Data analysis I: CRP mixture regression (bulk-only) with model-based clustering}
\label{sec:app-clustering}

This application demonstrates model-based clustering from a
\emph{conditional} DPM regression, using the
CRP backend and the bulk-only likelihood
(no GPD tail). Clusters are inferred from posterior draws of the latent
membership indicators and summarized using a posterior similarity matrix
(PSM), which is invariant to label switching.

\subsection{Dataset and objective}

I use the \code{Boston} housing data (\pkg{MASS}), treating each census
tract as an observational unit. Let $$
  \mathcal{D}=\{(y_i,\boldsymbol{x}_i)\}_{i=1}^n,
$$ where $y_i$ is the per-capita crime rate \code{crim} for tract $i$
and $\boldsymbol{x}_i$ is a vector of tract characteristics. The goal is
twofold: (i) flexibly model the conditional distribution of a positive,
right-skewed outcome and (ii) identify latent groups of tracts with
distinct regression patterns.

\subsection{CRP mixture regression model (bulk-only)}

I fit a DPM of lognormal regressions. Conditional on a latent
membership $z_i\in\{1,2,\ldots\}$, the outcome follows the bulk kernel
introduced earlier: \begin{align}
  y_i \mid z_i=k,\boldsymbol{x}_i
  &\sim \mathrm{Lognormal}\!\left(\mu_{ik},\sigma_k\right), \label{eq:app_clust_kernel}\\
  \mu_{ik}
  &= g\!\left(\boldsymbol{x}_i^\top \boldsymbol{\beta}_k\right),
  \qquad g(\eta)=\eta, \label{eq:app_clust_link}
\end{align} where $\boldsymbol{\beta}_k$ are component-specific
regression coefficients and $\sigma_k>0$ is the component-specific
log-scale standard deviation. The DP prior induces a random partition
through the CRP representation, $$
  (z_1,\ldots,z_n)\mid \alpha \sim \mathrm{CRP}(\alpha),
$$ with $\alpha>0$ controlling the expected number of occupied clusters.
In \pkg{CausalMixGPD}, the CRP backend represents up to
\code{components = K} clusters in a finite NIMBLE model; the realized
number of occupied clusters is then monitored and checked against this
upper bound.

\subsection{Posterior simulation, convergence checks, and parameter extraction}

Posterior simulation is carried out by MCMC. I run two chains and
assess convergence using standard diagnostics (trace plots,
autocorrelation, effective sample sizes, and $\widehat{R}$ when
applicable), along with a mixture-specific diagnostic: the posterior
distribution of the number of occupied clusters.

A key practical point for this analysis is that we explicitly request
link-mode for \code{meanlog} so the CRP model is a
\emph{mixture regression} rather than a purely unconditional mixture.
This is done via \code{param\_specs} in the wrapper call below.



\begin{knitrout}\footnotesize
\definecolor{shadecolor}{rgb}{0.969, 0.969, 0.969}\color{fgcolor}\begin{kframe}
\begin{alltt}
\hldef{fit_clust} \hlkwb{<-} \hlkwd{dpmix}\hldef{(}
  \hlkwc{formula} \hldef{= crim} \hlopt{~} \hldef{lstat} \hlopt{+} \hldef{rm} \hlopt{+} \hldef{nox,}
  \hlkwc{data} \hldef{= train_dat,}
  \hlkwc{backend} \hldef{=} \hlsng{"crp"}\hldef{,}
  \hlkwc{kernel} \hldef{=} \hlsng{"normal"}\hldef{,}
  \hlkwc{components} \hldef{=} \hlnum{5}\hldef{,}
  \hlkwc{mcmc} \hldef{= mcmc_fixed}
\hldef{)}
\end{alltt}
\end{kframe}
\end{knitrout}

\begin{knitrout}\footnotesize
\definecolor{shadecolor}{rgb}{0.969, 0.969, 0.969}\color{fgcolor}\begin{kframe}
\begin{verbatim}
  sample   n
1  train 404
2    new 102
\end{verbatim}
\end{kframe}
\end{knitrout}

\subsection{Allocation plot: training vs new data}

For clustering, we report old (training) and new-sample cluster tables
and a bar chart comparing cluster frequencies.



\begin{knitrout}\footnotesize
\definecolor{shadecolor}{rgb}{0.969, 0.969, 0.969}\color{fgcolor}\begin{kframe}
\begin{alltt}
\hldef{cluster_alloc_new} \hlkwb{<-} \hlkwd{allocation}\hldef{(fit_clust,} \hlkwc{newdata} \hldef{= newdata_cluster)}
\end{alltt}
\end{kframe}
\end{knitrout}

\begin{knitrout}\footnotesize
\definecolor{shadecolor}{rgb}{0.969, 0.969, 0.969}\color{fgcolor}\begin{kframe}
\begin{verbatim}
  cluster n_train
1       1     238
2       2      50
3       3      19
4       4      82
5       5      15
  cluster n_new
1       1    55
2       2    12
3       3    13
4       4    18
5       5     4
  sample       min       q25    median       q75       max
1    old 0.3595995 0.5446042 0.8469694 0.8742747 0.8954715
2    new 0.2732306 0.5608872 0.8181139 0.9030104 0.9934707
\end{verbatim}
\end{kframe}
\end{knitrout}

\begin{knitrout}\footnotesize
\definecolor{shadecolor}{rgb}{0.969, 0.969, 0.969}\color{fgcolor}\begin{kframe}
\begin{alltt}
\hlkwd{ggplot}\hldef{(cluster_old_new_counts,} \hlkwd{aes}\hldef{(}\hlkwc{x} \hldef{= cluster,} \hlkwc{y} \hldef{= n,} \hlkwc{fill} \hldef{= sample))} \hlopt{+}
  \hlkwd{geom_col}\hldef{(}\hlkwc{position} \hldef{=} \hlsng{"dodge"}\hldef{)} \hlopt{+}
  \hlkwd{labs}\hldef{(}
    \hlkwc{title} \hldef{=} \hlsng{"Cluster frequencies: old vs new"}\hldef{,}
    \hlkwc{x} \hldef{=} \hlsng{"Cluster"}\hldef{,}
    \hlkwc{y} \hldef{=} \hlsng{"Count"}\hldef{,}
    \hlkwc{fill} \hldef{=} \hlsng{"Sample"}
  \hldef{)} \hlopt{+}
  \hlkwd{theme_minimal}\hldef{(}\hlkwc{base_size} \hldef{=} \hlnum{10}\hldef{)}
\end{alltt}
\end{kframe}
\end{knitrout}

\begin{knitrout}\footnotesize
\definecolor{shadecolor}{rgb}{0.969, 0.969, 0.969}\color{fgcolor}

{\centering \includegraphics[width=0.82\linewidth]{cache/figs/app_cluster_old_new_barplot} 

}


\end{knitrout}

\section{Data analysis II: Spliced causal mixture regression with ATE, QTE, CATE and CQTE}
\label{sec:app-causal}

This application demonstrates causal estimands from a two-arm spliced
DPM regression with a GPD tail. In contrast to Section~\ref{sec:app-clustering}, where the
focus is model-based clustering under a CRP bulk-only regression, here
the focus is inference on mean and quantile treatment contrasts,
including conditional (covariate-specific) effects and extreme
upper-tail behavior.

\subsection{Dataset and objective}

The \pkg{AER} package does not ship the National JTPA experimental data
as an in-package dataset. Instead, we use a bundled causal benchmark
data object to provide a clean and reproducible entry point for a
heavy-tailed earnings-style application.

Let $Y$ denote post-program earnings (USD), $A\in\{0,1\}$ denote
randomized assignment to training, and $\boldsymbol{x}$ denote baseline
covariates including previous earnings and demographics. The scientific
goal is to quantify how training changes (i) average earnings and (ii)
distributional features of earnings, with special attention to the upper
tail.

Throughout, we report four causal estimands: \begin{align*}
  \mathrm{ATE} &= \mathbb{E}\{Y(1)\}-\mathbb{E}\{Y(0)\},\\
  \mathrm{QTE}(\tau) &= Q_{Y(1)}(\tau)-Q_{Y(0)}(\tau),\\
  \mathrm{CATE}(\boldsymbol{x}) &= \mathbb{E}\{Y(1)\mid \boldsymbol{X}=\boldsymbol{x}\}-\mathbb{E}\{Y(0)\mid \boldsymbol{X}=\boldsymbol{x}\},\\
  \mathrm{CQTE}(\tau,\boldsymbol{x}) &= Q_{Y(1)\mid \boldsymbol{X}=\boldsymbol{x}}(\tau)-Q_{Y(0)\mid \boldsymbol{X}=\boldsymbol{x}}(\tau),
\end{align*} where $Q(\tau)$ denotes the $\tau$-quantile. The
standardized estimands (ATE, QTE) integrate over the empirical
distribution of $\{\boldsymbol{x}_i\}_{i=1}^n$; the conditional
estimands (CATE, CQTE) are evaluated at user-chosen covariate profiles.



\subsection{Spliced DPM regression model}

I model arm-specific conditional outcome distributions
$F_a(\cdot\mid \boldsymbol{x})$ via a spliced construction: $$
  f_a(y\mid \boldsymbol{x})
  =
  \sum_j w_{aj}(\boldsymbol{x})\, k(y\mid \Theta_{aj})\,\mathbf{1}(y\le u)
  +
  \{1-F_a(u\mid \boldsymbol{x})\}\, g(y\mid u,\sigma_a,\xi_a)\,\mathbf{1}(y>u),
$$ where the bulk kernel $k(\cdot)$ is a DPM regression
(stick-breaking backend) and the tail density $g(\cdot)$ is GPD with
scale $\sigma_a>0$ and shape $\xi_a$. The GPD tail index $\xi_a$
controls extremal heaviness (e.g., $\xi_a>0$ corresponds to heavy
tails). This model stabilizes inference for extreme quantiles by
coupling nonparametric bulk flexibility with EVT-consistent tail
extrapolation.

\subsection{Model fitting}

I fit a two-arm spliced mixture using a lognormal bulk kernel. The MCMC
settings below are intentionally modest for vignette speed; in
substantive analysis we recommend longer runs and sensitivity checks
(e.g., components, threshold priors, and prior scales).



\begin{knitrout}\footnotesize
\definecolor{shadecolor}{rgb}{0.969, 0.969, 0.969}\color{fgcolor}\begin{kframe}
\begin{alltt}
\hldef{mcmc_fixed} \hlkwb{<-} \hlkwd{list}\hldef{(}
  \hlkwc{niter} \hldef{=} \hlnum{120}\hldef{,}
  \hlkwc{nburnin} \hldef{=} \hlnum{40}\hldef{,}
  \hlkwc{thin} \hldef{=} \hlnum{1}\hldef{,}
  \hlkwc{nchains} \hldef{=} \hlnum{1}\hldef{,}
  \hlkwc{seed} \hldef{=} \hlnum{101}\hldef{,}
  \hlkwc{show_progress} \hldef{=} \hlnum{FALSE}
\hldef{)}
\end{alltt}
\end{kframe}
\end{knitrout}

\begin{knitrout}\footnotesize
\definecolor{shadecolor}{rgb}{0.969, 0.969, 0.969}\color{fgcolor}\begin{kframe}
\begin{alltt}
\hldef{fit} \hlkwb{<-} \hlkwd{dpmgpd}\hldef{(}
  \hlkwc{x} \hldef{= y,}
  \hlkwc{X} \hldef{= X,}
  \hlkwc{treat} \hldef{= A,}
  \hlkwc{backend} \hldef{=} \hlsng{"sb"}\hldef{,}
  \hlkwc{kernel} \hldef{=} \hlsng{"normal"}\hldef{,}
  \hlkwc{components} \hldef{=} \hlnum{4}\hldef{,}
  \hlkwc{mcmc} \hldef{= mcmc_fixed}
\hldef{)}
\end{alltt}
\end{kframe}
\end{knitrout}

\begin{knitrout}\footnotesize
\definecolor{shadecolor}{rgb}{0.969, 0.969, 0.969}\color{fgcolor}\begin{kframe}
\begin{verbatim}
     metric value
1         n   500
2         p     3
3 n_treated   309
4 n_control   191
\end{verbatim}
\end{kframe}
\end{knitrout}

\subsection{Standardized estimands: ATE and QTE}

I first summarize the overall average effect using \texttt{ate()},
which aggregates draw-wise conditional mean contrasts over the empirical
covariate distribution. I then compute quantile effects at a grid
including upper-tail levels.



\begin{knitrout}\footnotesize
\definecolor{shadecolor}{rgb}{0.969, 0.969, 0.969}\color{fgcolor}\begin{kframe}
\begin{alltt}
\hldef{ate_fit} \hlkwb{<-} \hlkwd{ate}\hldef{(fit,} \hlkwc{show_progress} \hldef{=} \hlnum{FALSE}\hldef{)}
\end{alltt}
\end{kframe}
\end{knitrout}

\begin{knitrout}\footnotesize
\definecolor{shadecolor}{rgb}{0.969, 0.969, 0.969}\color{fgcolor}\begin{kframe}
\begin{verbatim}
  id   estimate      lower      upper
1  1 -0.4121421 -0.7170418 0.09078951
\end{verbatim}
\end{kframe}
\end{knitrout}

\begin{knitrout}\footnotesize
\definecolor{shadecolor}{rgb}{0.969, 0.969, 0.969}\color{fgcolor}\begin{kframe}
\begin{alltt}
\hldef{qte_fit} \hlkwb{<-} \hlkwd{qte}\hldef{(fit,} \hlkwc{probs} \hldef{= taus,} \hlkwc{show_progress} \hldef{=} \hlnum{FALSE}\hldef{)}
\end{alltt}
\end{kframe}
\end{knitrout}

\begin{knitrout}\footnotesize
\definecolor{shadecolor}{rgb}{0.969, 0.969, 0.969}\color{fgcolor}\begin{kframe}
\begin{verbatim}
  id index   estimate       lower      upper
1  1  0.25 -0.8197706 -1.48164947 -0.2415758
2  1  0.50 -0.6946461 -1.51414175 -0.1011650
3  1  0.75  0.7299151 -0.25027110  1.2937407
4  1  0.90  0.7794739 -0.08451773  1.6423832
5  1  0.95  0.7929622 -0.89400954  2.1972027
6  1  0.99  1.0779570 -5.36622364  5.7102460
\end{verbatim}
\end{kframe}
\end{knitrout}

\begin{knitrout}\footnotesize
\definecolor{shadecolor}{rgb}{0.969, 0.969, 0.969}\color{fgcolor}\begin{kframe}
\begin{alltt}
\hldef{ate_effect_plot} \hlkwb{<-} \hlkwd{plot}\hldef{(ate_fit,} \hlkwc{type} \hldef{=} \hlsng{"effect"}\hldef{)}
\hldef{ate_effect_plot}
\end{alltt}
\end{kframe}
\end{knitrout}

\begin{knitrout}\footnotesize
\definecolor{shadecolor}{rgb}{0.969, 0.969, 0.969}\color{fgcolor}\begin{kframe}
\begin{alltt}
\hldef{qte_effect_plot} \hlkwb{<-} \hlkwd{plot}\hldef{(qte_fit,} \hlkwc{type} \hldef{=} \hlsng{"effect"}\hldef{)}
\hldef{qte_effect_plot}
\end{alltt}
\end{kframe}
\end{knitrout}

\begin{knitrout}\footnotesize
\definecolor{shadecolor}{rgb}{0.969, 0.969, 0.969}\color{fgcolor}

{\centering \includegraphics[width=0.82\linewidth]{cache/figs/app2_ate_effect} 

}


\end{knitrout}

\begin{knitrout}\footnotesize
\definecolor{shadecolor}{rgb}{0.969, 0.969, 0.969}\color{fgcolor}

{\centering \includegraphics[width=0.82\linewidth]{cache/figs/app2_qte_effect} 

}


\end{knitrout}

\subsection{Conditional estimands: CATE and CQTE at representative profiles}

To illustrate heterogeneity, we evaluate conditional effects at two
covariate profiles: (i) a
`typical'' profile at sample medians and (ii) a`high previous earnings''
profile at the 80th percentile of previous earnings, holding other
covariates at medians.



\begin{knitrout}\footnotesize
\definecolor{shadecolor}{rgb}{0.969, 0.969, 0.969}\color{fgcolor}\begin{kframe}
\begin{alltt}
\hldef{col_med} \hlkwb{<-} \hlkwd{apply}\hldef{(X,} \hlnum{2}\hldef{, median,} \hlkwc{na.rm} \hldef{=} \hlnum{TRUE}\hldef{)}
\hldef{x_typical} \hlkwb{<-} \hldef{col_med}
\hldef{x_highprev} \hlkwb{<-} \hldef{col_med}
\hldef{profile_var} \hlkwb{<-} \hlkwa{if} \hldef{(}\hlsng{"prevearn"} \hlopt{%in%} \hlkwd{colnames}\hldef{(X))} \hlsng{"prevearn"} \hlkwa{else} \hlkwd{colnames}\hldef{(X)[}\hlnum{1}\hldef{]}
\hldef{x_highprev[profile_var]} \hlkwb{<-} \hlkwd{as.numeric}\hldef{(}\hlkwd{quantile}\hldef{(X[, profile_var],} \hlnum{0.80}\hldef{,} \hlkwc{na.rm} \hldef{=} \hlnum{TRUE}\hldef{))}
\hldef{xnew} \hlkwb{<-} \hlkwd{rbind}\hldef{(}\hlkwc{typical} \hldef{= x_typical,} \hlkwc{high_prev_earn} \hldef{= x_highprev)}
\end{alltt}
\end{kframe}
\end{knitrout}

\begin{knitrout}\footnotesize
\definecolor{shadecolor}{rgb}{0.969, 0.969, 0.969}\color{fgcolor}\begin{kframe}
\begin{verbatim}
  profile_var
1          x1
                       x1          x2         x3
typical        0.07649836 -0.09628321 0.01209002
high_prev_earn 0.79678466 -0.09628321 0.01209002
\end{verbatim}
\end{kframe}
\end{knitrout}



\begin{knitrout}\footnotesize
\definecolor{shadecolor}{rgb}{0.969, 0.969, 0.969}\color{fgcolor}\begin{kframe}
\begin{alltt}
\hldef{cate_fit} \hlkwb{<-} \hlkwd{cate}\hldef{(fit,} \hlkwc{newdata} \hldef{= xnew,} \hlkwc{show_progress} \hldef{=} \hlnum{FALSE}\hldef{)}
\end{alltt}
\end{kframe}
\end{knitrout}

\begin{knitrout}\footnotesize
\definecolor{shadecolor}{rgb}{0.969, 0.969, 0.969}\color{fgcolor}\begin{kframe}
\begin{verbatim}
  id   estimate      lower      upper
1  1 -0.3645759 -0.6971991  0.1087821
2  2 -1.4191524 -2.2463081 -0.6754342
\end{verbatim}
\end{kframe}
\end{knitrout}

\begin{knitrout}\footnotesize
\definecolor{shadecolor}{rgb}{0.969, 0.969, 0.969}\color{fgcolor}\begin{kframe}
\begin{alltt}
\hldef{cqte_fit} \hlkwb{<-} \hlkwd{cqte}\hldef{(fit,} \hlkwc{probs} \hldef{= taus,} \hlkwc{newdata} \hldef{= xnew,} \hlkwc{show_progress} \hldef{=} \hlnum{FALSE}\hldef{)}
\end{alltt}
\end{kframe}
\end{knitrout}

\begin{knitrout}\footnotesize
\definecolor{shadecolor}{rgb}{0.969, 0.969, 0.969}\color{fgcolor}\begin{kframe}
\begin{verbatim}
   id index   estimate      lower       upper
1   1  0.25 -0.2157592 -0.5805429  0.08519506
2   1  0.50 -0.2653653 -0.4044642 -0.07230213
3   1  0.75 -0.2282270 -0.4720426  0.05270702
4   1  0.90 -0.1783537 -0.8925062  0.33589134
5   1  0.95 -0.1166213 -1.2130864  0.59538829
6   1  0.99 -0.6052849 -3.2226522  0.81979673
7   2  0.25 -1.0029540 -1.4523082 -0.60897787
8   2  0.50 -0.7302790 -1.2913572 -0.17271939
9   2  0.75 -1.1369796 -2.1213103 -0.57987170
10  2  0.90 -2.5609432 -3.9700641 -1.28956486
11  2  0.95 -2.5888083 -4.2046625 -1.28730194
12  2  0.99 -2.7664670 -5.7033867 -0.92284556
\end{verbatim}
\end{kframe}
\end{knitrout}

CATE/CQTE effect plots are omitted in this PDF.

\subsection{Tail behavior}

A key advantage of the spliced model is direct posterior inference on
the tail index $\xi_a$ by arm. Larger $\xi_a$ indicates a heavier upper
tail and greater propensity for extreme earnings.



\begin{knitrout}\footnotesize
\definecolor{shadecolor}{rgb}{0.969, 0.969, 0.969}\color{fgcolor}\begin{kframe}
\begin{verbatim}
        arm      mean         sd    q0.025    q0.500
23  treated 0.5830602 0.04488664 0.4780895 0.5805598
231 control 0.3667860 0.11154470 0.1951404 0.4457394
       q0.975
23  0.6883067
231 0.4635703
\end{verbatim}
\end{kframe}
\end{knitrout}

\subsection{Interpretation}

The ATE summarizes the average shift in earnings under training, while
QTE describes how the effect varies across the marginal outcome
distribution. CATE and CQTE extend these contrasts to covariate-specific
comparisons, illustrating treatment effect heterogeneity in both central
and extreme quantiles. Finally, arm-specific tail indices $\xi_a$
provide a parsimonious summary of whether training changes extremal
behavior beyond what is captured by bulk-only regression.


\section{Discussion}
\label{sec:discussion}

\section{Supplementary Materials}
\label{sec:supplementary}

\bibliography{CausalMixGPD_JSS_refs}

\end{document}


